%%% The main file. It contains definitions of basic parameters and includes all other parts.

%% Settings for single-side (simplex) printing
% Margins: left 40mm, right 25mm, top and bottom 25mm
% (but beware, LaTeX adds 1in implicitly)
\documentclass[12pt,a4paper]{report}
\setlength\textwidth{145mm}
\setlength\textheight{247mm}
\setlength\oddsidemargin{15mm}
\setlength\evensidemargin{15mm}
\setlength\topmargin{0mm}
\setlength\headsep{0mm}
\setlength\headheight{0mm}
% \openright makes the following text appear on a right-hand page
\let\openright=\clearpage

%% Settings for two-sided (duplex) printing
% \documentclass[12pt,a4paper,twoside,openright]{report}
% \setlength\textwidth{145mm}
% \setlength\textheight{247mm}
% \setlength\oddsidemargin{14.2mm}
% \setlength\evensidemargin{0mm}
% \setlength\topmargin{0mm}
% \setlength\headsep{0mm}
% \setlength\headheight{0mm}
% \let\openright=\cleardoublepage

%% Generate PDF/A-2u
\usepackage[a-2u]{pdfx}

%% Character encoding: usually latin2, cp1250 or utf8:
\usepackage[utf8]{inputenc}

%% Prefer Latin Modern fonts
\usepackage{lmodern}

%% Further useful packages (included in most LaTeX distributions)
\usepackage{amsmath}        % extensions for typesetting of math
\usepackage{amsfonts}       % math fonts
\usepackage{amsthm}         % theorems, definitions, etc.
\usepackage{amssymb}
\usepackage{indentfirst}
\usepackage{bbding}         % various symbols (squares, asterisks, scissors, ...)
\usepackage{bm}             % boldface symbols (\bm)
\usepackage{graphicx}       % embedding of pictures
\usepackage{fancyvrb}       % improved verbatim environment
\usepackage{natbib}         % citation style AUTHOR (YEAR), or AUTHOR [NUMBER]
\usepackage[nottoc]{tocbibind} % makes sure that bibliography and the lists
			    % of figures/tables are included in the table
			    % of contents
\usepackage{dcolumn}        % improved alignment of table columns
\usepackage{booktabs}       % improved horizontal lines in tables
\usepackage{paralist}       % improved enumerate and itemize
\usepackage{xcolor}         % typesetting in color

\bibliographystyle{alpha}
%%% My packages

\usepackage{stmaryrd}
\usepackage{mathtools}

%%% Basic information on the thesis

% Thesis title in English (exactly as in the formal assignment)
\def\ThesisTitle{Pseudofinite structures and limits}

% Author of the thesis
\def\ThesisAuthor{Ondřej Ježil}

% Year when the thesis is submitted
\def\YearSubmitted{2022}

% Name of the department or institute, where the work was officially assigned
% (according to the Organizational Structure of MFF UK in English,
% or a full name of a department outside MFF)
\def\Department{Department of Algebra}

% Is it a department (katedra), or an institute (ústav)?
\def\DeptType{Department}

% Thesis supervisor: name, surname and titles
\def\Supervisor{prof. RNDr. Jan Krajíček Ph.D.}

% Supervisor's department (again according to Organizational structure of MFF)
\def\SupervisorsDepartment{Department of Algebra}

% Study programme and specialization
\def\StudyProgramme{Mathematics}
\def\StudyBranch{Mathematics for Information Technologies}

% An optional dedication: you can thank whomever you wish (your supervisor,
% consultant, a person who lent the software, etc.)
\def\Dedication{%
I dedicate this thesis to my parents. I also want to thank my supervisor, professor \text{Jan Krajíček}, for both his expert knowledge and quick attendance to my numerous questions.
}

% Abstract (recommended length around 80-200 words; this is not a copy of your thesis assignment!)
\def\Abstract{%
For a class of graph instances of a computational problem we define a limit object, relative to some computationally restricted class of functions. The key method here is forcing with random variables where the sample set is taken as instances of some nonstandard size. We study the general theory of these limits, called in the thesis wide limits, and their connection to classical problems such as finding a large clique or with the combinatorial problems associated with the classes of total $\NP$ search problems $\PPA$ and $\PPAD$. Our main results are several 0-1 laws associated with these limits and existence of a significantly large clique of the wide limit of all graph consisting of one large clique.
}

% 3 to 5 keywords (recommended), each enclosed in curly braces
\def\Keywords{%
{pseudofinite structure}, {forcing}, {complexity}, {witnessing}, {\TFNP}, {clique}
}

%% The hyperref package for clickable links in PDF and also for storing
%% metadata to PDF (including the table of contents).
%% Most settings are pre-set by the pdfx package.
\hypersetup{unicode}
\hypersetup{breaklinks=true}

% Definitions of macros (see description inside)
%%% This file contains definitions of various useful macros and environments %%%
%%% Please add more macros here instead of cluttering other files with them. %%%

%%% Minor tweaks of style

% These macros employ a little dirty trick to convince LaTeX to typeset
% chapter headings sanely, without lots of empty space above them.
% Feel free to ignore.
\makeatletter
\def\@makechapterhead#1{
  {\parindent \z@ \raggedright \normalfont
   \Huge\bfseries \thechapter. #1
   \par\nobreak
   \vskip 20\p@
}}
\def\@makeschapterhead#1{
  {\parindent \z@ \raggedright \normalfont
   \Huge\bfseries #1
   \par\nobreak
   \vskip 20\p@
}}
\makeatother

% This macro defines a chapter, which is not numbered, but is included
% in the table of contents.
\def\chapwithtoc#1{
\chapter*{#1}
\addcontentsline{toc}{chapter}{#1}
}

% Draw black "slugs" whenever a line overflows, so that we can spot it easily.
\overfullrule=1mm

%%% Macros for definitions, theorems, claims, examples, ... (requires amsthm package)

\theoremstyle{plain}
\newtheorem{thm}{Theorem}
\newtheorem{lemma}[thm]{Lemma}
\newtheorem{claim}[thm]{Claim}

\theoremstyle{plain}
\newtheorem{defn}{Definition}

\theoremstyle{remark}
\newtheorem*{cor}{Corollary}
\newtheorem*{rem}{Remark}
\newtheorem*{example}{Example}

%%% An environment for proofs

\newenvironment{myproof}{
  \par\medskip\noindent
  \textit{Proof}.
}{
\newline
\rightline{$\qedsymbol$}
}

%%% An environment for typesetting of program code and input/output
%%% of programs. (Requires the fancyvrb package -- fancy verbatim.)

\DefineVerbatimEnvironment{code}{Verbatim}{fontsize=\small, frame=single}

%%% The field of all real and natural numbers
\newcommand{\R}{\mathbb{R}}
\newcommand{\N}{\mathbb{N}}

%%% Useful operators for statistics and probability
\DeclareMathOperator{\pr}{\textsf{P}}
\DeclareMathOperator{\E}{\textsf{E}\,}
\DeclareMathOperator{\var}{\textrm{var}}
\DeclareMathOperator{\sd}{\textrm{sd}}

%%% Transposition of a vector/matrix
\newcommand{\T}[1]{#1^\top}

%%% Various math goodies
\newcommand{\goto}{\rightarrow}
\newcommand{\gotop}{\stackrel{P}{\longrightarrow}}
\newcommand{\maon}[1]{o(n^{#1})}
\newcommand{\abs}[1]{\left|{#1}\right|}
\newcommand{\dint}{\int_0^\tau\!\!\int_0^\tau}
\newcommand{\isqr}[1]{\frac{1}{\sqrt{#1}}}

%%% Various table goodies
\newcommand{\pulrad}[1]{\raisebox{1.5ex}[0pt]{#1}}
\newcommand{\mc}[1]{\multicolumn{1}{c}{#1}}

%%% Custom macros
\newcommand{\Th}{\text{Th}}
\newcommand{\st}{\text{st}}
\newcommand{\bbl}{\llbracket}
\newcommand{\bbr}{\rrbracket}
\newcommand{\Lor}{\bigvee}
\newcommand{\Land}{\bigwedge}
\newcommand{\1}{\textbf{1}}
\newcommand{\0}{\textbf{0}}
\newcommand{\dpt}{\text{dp}}
\newcommand{\Aut}{\text{Aut}}
\newcommand{\pPATH}{*\text{PATH}}
\newcommand{\EDGE}{\text{EDGE}}
\newcommand{\nEDGE}{\text{nonEDGE}}
\newcommand{\ALL}{\text{ALL}}
\newcommand{\leqp}{\trianglelefteq}
\newcommand{\set}[1]{\langle #1)}
\newcommand{\ceil}[1]{\left \lceil {#1} \right \rceil}
\newcommand{\floor}[1]{\left \lfloor {#1} \right \rfloor}
\newcommand{\M}{\mathcal{M}}
\newcommand{\Po}{\mathcal{P}}
\newcommand{\A}{\mathcal{A}}
\newcommand{\B}{\mathcal{B}}
\newcommand{\I}{\mathcal{I}}
\newcommand{\card}{\text{card}}
\newcommand{\SK}{\text{SK}}
\newcommand{\CK}{\text{CK}}
\newcommand{\G}{\mathcal{G}}
\newcommand{\Q}{\mathbb{Q}}
\renewcommand{\P}{\textbf{P}}
\newcommand{\NP}{\textbf{NP}}
\newcommand{\FP}{\textbf{FP}}
\newcommand{\TFNP}{\textbf{TFNP}}
\newcommand{\LEAF}{\text{LEAF}}
\newcommand{\SOURCEORSINK}{\text{SOURCE.OR.SINK}}
\newcommand{\PPA}{\textbf{PPA}}
\newcommand{\PPAD}{\textbf{PPAD}}
\renewcommand{\N}{\mathbb{N}}
\newcommand{\dom}{\text{dom}\:}
\newcommand{\Rng}{\text{Rng}\:}
\newcommand{\Frac}{\text{Frac}}
%\newcommand{\abs}[1]{\lvert #1 \rvert}


% Definitions of amsthm environments
\newtheoremstyle{mystyle}%                % Name
  {}%                                     % Space above
  {}%                                     % Space below
  {}%                             % Body font
  {}%                                     % Indent amount
  {\bfseries}%                            % Theorem head font
  {.}%                                    % Punctuation after theorem head
  { }%                                    % Space after theorem head, ' ', or \newline
  {}%                                     % Theorem head spec (can be left empty, meaning `normal')

\theoremstyle{mystyle}

\newtheorem{thrm}{Theorem}[section]
\newtheorem{lemm}[thrm]{Lemma}
\newtheorem{conj}[thrm]{Conjecture}
\newtheorem{ques}[thrm]{Question}
\newtheorem*{prope}{Property}
\newtheorem*{thrm*}{Theorem}
\newtheorem{exam}[thrm]{Example}
\newtheorem{defi}[thrm]{Definition}
\newtheorem{crll}[thrm]{Corollary}

% Title page and various mandatory informational pages
\begin{document}
%%% Title page of the thesis and other mandatory pages

%%% Title page of the thesis

\pagestyle{empty}
\hypersetup{pageanchor=false}
\begin{center}

\centerline{\mbox{\includegraphics[width=166mm]{../img/logo-en.pdf}}}

\vspace{-8mm}
\vfill

{\bf\Large MASTER THESIS}

\vfill

{\LARGE\ThesisAuthor}

\vspace{15mm}

{\LARGE\bfseries\ThesisTitle}\footnote{This is a second version of a thesis originally defended in June 2022. It includes errata and several parts were rewritten to improve clarity.}

\vfill

\Department

\vfill

{
\centerline{\vbox{\halign{\hbox to 0.45\hsize{\hfil #}&\hskip 0.5em\parbox[t]{0.45\hsize}{\raggedright #}\cr
Supervisor of the master thesis:&\Supervisor \cr
\noalign{\vspace{2mm}}
Study programme:&\StudyProgramme \cr
\noalign{\vspace{2mm}}
Study branch:&\StudyBranch \cr
}}}}

\vfill

% Zde doplňte rok
Prague \YearSubmitted

\end{center}

\newpage

%%% Here should be a bound sheet included -- a signed copy of the "master
%%% thesis assignment". This assignment is NOT a part of the electronic
%%% version of the thesis. DO NOT SCAN.

%%% A page with a solemn declaration to the master thesis

%%%\openright
%%%\hypersetup{pageanchor=true}
%%%\pagestyle{plain}
%%%\pagenumbering{roman}
%%%\vglue 0pt plus 1fill
%%%
%%%\noindent
%%%I declare that I carried out this master thesis independently, and only with the cited
%%%sources, literature and other professional sources. It has not been used to obtain another
%%%or the same degree.
%%%
%%%\medskip\noindent
%%%I understand that my work relates to the rights and obligations under the Act No.~121/2000 Sb.,
%%%the Copyright Act, as amended, in particular the fact that the Charles
%%%University has the right to conclude a license agreement on the use of this
%%%work as a school work pursuant to Section 60 subsection 1 of the Copyright~Act.
%%%
%%%\vspace{10mm}
%%%
%%%\hbox{\hbox to 0.5\hsize{%
%%%In \hbox to 6em{\dotfill} date \hbox to 6em{\dotfill}
%%%\hss}\hbox to 0.5\hsize{\dotfill\quad}}
%%%\smallskip
%%%\hbox{\hbox to 0.5\hsize{}\hbox to 0.5\hsize{\hfil Author's signature\hfil}}
%%%
%%%\vspace{20mm}
%%%\newpage
%%%
%%%%%% Dedication
%%%
%%%\openright
%%%
%%%\noindent
%%%\Dedication
%%%
%%%\newpage

%%% Mandatory information page of the thesis

\openright

\vbox to 0.5\vsize{
\setlength\parindent{0mm}
\setlength\parskip{5mm}

Title:
\ThesisTitle

Author:
\ThesisAuthor

\DeptType:
\Department

Supervisor:
\Supervisor, \SupervisorsDepartment

Abstract:
\Abstract

Keywords:
\Keywords

\vss}

\newpage

\openright
\pagestyle{plain}
\pagenumbering{arabic}
\setcounter{page}{1}


%%% A page with automatically generated table of contents of the master thesis

\tableofcontents

%%% Each chapter is kept in a separate file
\chapter*{Introduction}
\addcontentsline{toc}{chapter}{Introduction}

There exist several logical constructions of limits of classes of finite structures such as the ultraproduct construction or using the compactness theorem. The latter used in \cite{Fagin1976} to prove the 0-1 law for structures over relational vocabularies.

In combinatorics there are also several notions of limits of finite graphs. For example the dense graph limit defined for a sequence of graph $\{G_k\}_{k>0}$ satisfying the condition that
\[t(F,G_n)=\frac{\hom(F,G)}{\abs{G_n}^{\abs{F}}}\]
converges for every fixed connected graph $F$ which provided a framework (see \cite{lovasz2006limits}) to restate and find new proofs for results in extremal graph theory such as Goodman's theorem relating the number of edges to the number of triangles in a graph. There are other notions of limits of sequences of graphs we recommend to read \cite{Nesetril2013} to the interested reader. Another recent use of limit objects for the results of extremal combinatorics was by Razborov in \cite{razborov2007flag}.

These diferent notions of limits directly or tangentially relate to the concept of pseudofinite structures. For a first order language $L$, we call an $L$-structure $S$ pseudofinite, if it satisfies the theory $T_f$ consisting of all sentences true in all finite $L$-structures.

In this thesis we use the concept of pseudofinite structures to define a limit of a family of finite graphs relative to some computationally restricted class of functions. Instead of studying the density of substructures we study these so called wide limits both generally and by analyzing concrete examples and tying them with complexity-theoretical statements.

The key method we use is arithmetical forcing with random variables which allows us to construct models of (weak) arithmetical theories and by restricting to a language of graphs gives us boolean valued graphs. An important resource for this technique is \cite{krajicek2010forcing}. In these boolean valued graphs, witnessing of existential quantifiers corresponds to complexity theoretical properties of the class of graphs we are considering. 

(TODO: Comment on structure of chapters.)

\chapter*{Preliminaries}
\addcontentsline{toc}{chapter}{Preliminaries}

In this chapter we recall a few important notions which we will use in the next chapter to define the central construction we study. We do not review the notions formally but always provide a reference for the reader unfamiliar with these topics. Throughout this thesis we assume a basic knowledge of mathematical logic, model theory and measure theory. Important concepts for us are the nonstandard models of true arithmetic, nonstandard analysis and $\NP$ search problems. We discuss these in more detail in the rest of this chapter.

\section*{The ambient model $\M$}
\addcontentsline{toc}{section}{The ambient model $\M$}

Let $L_{all}$ be the language consisting of function symbols for all functions on $\N$, all relations and all constants. We call a model of $\Th_{L_{all}}(\N)$ nonstandard if it is not isomorphic to $\N$. Every model of $\Th(\N)$ contains an initial segment isomorphic to $\N$ so we can view nonstandard models as those which also contain `infinite natural numbers', if we assume $\M\models \Th(\N)$ contains $\N$ as an initial segment then the elements $\M\setminus \N$ are called nonstandard. We recommend the introduction of \cite{kaye1991pa} for a review of this topic. In the appendix of \cite{krajicek2010forcing} there is an explict ultraproduct construction of a model $\M\models \Th(\N)$ which is $\aleph_1$-saturated.

This $\aleph_1$-saturated model $\M$ is used throughout this thesis and we call it the ambient model of arithmetic. For our applications we only need to know that the model is nonstandard and the following property holds because of the $\aleph_1$-saturation. (Note that we can encode finite sequences and sets in $\M$ which lets us state the property.)

\begin{prope}
If $\{a_k\}_{k\geq 0}$ is a sequence with elements in $\N$ then there is an element $t\in \M \setminus \N$ and a sequence $\{b_k\}_{k<t}\in \M$ with $a_k=b_k$ for all $k\in \N$.
\end{prope}

By overspill in $\M$ if some definable property $P$ holds for $a_k$ with unbounded indexes, then there is also some nonstandard $n<t$ such that $b_s$ satisfies the property $P$. Moreover, if some definable property $P$ holds for all $b_k$ with $k$ above some $k_0$ it has to also hold by induction in $\M$ for all nonstandard $b_n$. These $b_n$ are intuitively the limit elements of the sequence $\{a_k\}_{k\geq0}$.

Now let us introduce some notation. For $m\in\M$ we denote the set of numbers below $m$ as $\set{m}:=\{0,\dots,m-1\}\in\M$ and $\abs{m}$ as the bit-length of the (possibly-nonstandard) number $m$. These definitions can be easily made rigorous using the first order definitions of these functions.

\section*{Nonstandard analysis}
\addcontentsline{toc}{section}{Nonstandard analysis}

The reader can refer to \cite{goldbring2014lecture} for more formal treatment of topics discussed in this section including proofs. To use the method of forcing with random variables we need to consider the concept of $\M$-rationals. To define them we start by simply adjoining all negative elements to the semiring $\M$ to obtain the integral domain $\overline \M$. $\M$-rationals are then simply the ordered field of fractions $\Frac(\overline \M)$ which we denote $\Q^\M$.

There is a canonical injection $\Q \xhookrightarrow{} \Q^\M$ whose image consists exactly of the `standard fractions'. We call a $q\in\Q^\M$ finite if there is a standard $k$ such that $\abs{q}<\frac{k}{1}$, otherwise we call it infinite. We call $q\in\Q^\M$ infinitesimal if $q^{-1}$ is infinite (see also Definition \ref{defiinfinitesimal}). We will use $\Q^\M$ in a similar manner to what is called hyperreal numbers which lets us use it as an alternative foundation for the concepts of real analysis. The following is an important result which we use throughout the thesis.

\begin{thrm*}
There is a surjective function from the ring of finite $\M$-rationals to $\R$ which is a homomorphism of rings and the kernel is the ideal of infinitesimal numbers.

We denote the function $\st(-)$ and call it the \textbf{standard part of the $\M$-rational}.
\end{thrm*}

The following result characterizes convergence of a sequence of rational numbers in the language of nonstandard analysis.

\begin{thrm*}
Let $\{c_k\}_{k\geq 0}$ be a sequence of rational numbers and let $\{\tilde c_k\}_{t\geq k}$ be its nonstandard prolongation in $\M$. Then \[\lim_{k\to\infty}c_k = r \in \R\]
if and only if for any nonstandard $n\leq t$ we have that
$\st(\tilde c_n)=r.$
\end{thrm*}

We close this section with two inequalities heavily used in the proofs throughout the thesis.

\begin{thrm*}[Bernoulli's inequality]
Let $y \in \M$ and $x\in \Q^\M, x\geq -1$, then
\[(1+x)^y \geq 1+yx.\]
\end{thrm*}

\begin{thrm*}[Exponential inequality]
Let $y \in \M$, $x\in \Q^\M, x\geq 0$ and $x\leq y$, then
\[\left(1-\frac{x}{y}\right)^y \leq e^{-x}.\]
\end{thrm*}

\section*{Total $\NP$ search problems and polynomial oracle time}
\addcontentsline{toc}{section}{Total $\NP$ search problems and polynomial oracle time}

Our goal is to tie the properties of the wide limit with some complexity theoretic statements. We will recall several notions used later on.

The class of total $\NP$ search problems $\TFNP$, first defined in \cite{megiddo1991total}, consist of all relations on binary strings $P(x,y)$ such that: a) There is a polynomial time machine $M$ which, given $x,y$, can decide whether $P(x,y)$ holds. b) There exists a constant $c>0$ such that for every $x$ there exists at least one $y$ satisfying $\abs{y}\leq \abs{x}^c$ such that $P(x,y)$ holds.

While the definition of the class $\TFNP$ seems natural, the inner structure looks more arbitrary and the class is generally studied through its semantic subclasses. For example it is conjectured that there is no complete problem for $\TFNP$ \cite{goldberg2018tfnp} for a suitable notion of reducibility.

The semantic subclasses are defined as all problems reducible to some problem corresponding usually to a combinatorial lemma, for some appropriate definition of reduction. Two main subclasses are relevant for us. The class $\PPA$, polynomial parity argument, corresponds to all problems reducible to $\LEAF$, the problem formulated as follows. Given a number $k$ and a graph $G$ on the vertex set $\set{2^{\abs{k}}}$, presented by a boolean circuit of polynomial size in $\abs{k}$ computing its neighbourhood function, such that $\deg_G(0)=1$ and $\forall v: \deg_G(v)\leq 2$, find some nonzero $v$ with $\deg_G(v)=1$. The corresponding combinatorial principle being the handshaking lemma, which assures the problem is total. And the class $\PPAD$, directed polynomial parity argument, with the complete problem $\SOURCEORSINK$ formulated as follows. Given a number $k$ and a directed graph $G$ on the vertex set $\set{2^{\abs{k}}}$, presented by a boolean circuit of polynomial size in $\abs{k}$ computing the neighbourhood function, such that the vertex $0$ is a source and $\forall v: \deg_G^+(v),\deg_G^-(v)\leq 1$ find some nonzero vertex $v$ which is a source or a sink. The corresponding combinatorial principle here being the directed version of the handshaking lemma.

So far, we presented what is called `type 1' problem in \cite{beame1995rel}. The ones we are interested in are `type 2' problems which replace the input graph $G$ with a pair $(\alpha,x)$ consisting of an oracle $\alpha$ describing the neighbourhood function $N_G(-)$ (or both $N_G^+(-)$ and $N_G^-(-)$ in the directed case) on binary string of length at most $\abs{x}$. While the goal to solve these problems remains the same, suddenly the situation is quite different -- these type 2 classes can be easily separated from $\FP^2$, the type 2 version of polynomial time functions, by an adversary argument. More importantly we have that $\PPAD\subsetneq \PPA$ as type 2 complexity classes. Intuitively one can forget the orientation to get the undirected version, but cannot consistently assign orientation to undirected edges of a large graph.

The traditional model of computation is the Turing machines and for the type 2 problems oracle Turing machine. But to prove separations in the type 2 case, we can abstract the computation of an oracle Turing machine into a decision tree which describes the queries to an oracle. This abstraction is what to keep in mind in the following chapter when we define the class of functions $F_{rud}$.

\chapter{Forcing with random variables and the wide Limit}\label{chapwidelimit}
\section{Setup}

Our goal in this chapter is to provide a definition of a limit of a set an infinite set of finite graphs in which arbitrarily large graphs occur. The following definition makes our requirements of such a class of graphs precise. From now on $E$ denotes a binary relational symbol, so that formulas in the theory of graphs are precisely $E$-sentences.

\begin{defi}
Let $\{\mathcal{G}_k\}_{k>0}$ be a sequence of non-empty finite sets of (directed or undirected) finite graphs, i.e. structures of in the first order language consisting of the binary function $\{E\}$. We call it a \textbf{wide sequence} if the following holds.

\begin{itemize}
\item There is an increasing sequence of positive whole numbers $\{g_k\}_{k>0}$ such that the underlying set of each $G\in\mathcal{G}_k$ is $\set{g_k}$.
\item $\lim_{k\to\infty} \abs{\mathcal{G}_k}=\infty$
\end{itemize}
\end{defi}

The second condition guarantees that $\mathcal{G}_n$ is an infinite set for $n>\mathbb{N}$. In fact, the condition is not strictly important to proceed with this chapter, but it more closely describes what sort of sequences we are interested in. Many interesting classes of graphs form a wide sequence if we restrict the vertex-sets to $\set{g_k}$, where $\{g_k\}_{k>0}$ can be taken as the increasing sequence of all cardinalities in such a class. For example graphs with exactly one edge, graphs with bounded degree and so on. We will get to explore many examples in depth after we define the wide limit.

(TODO: wide sequence <-> class of structures)

 The limit is defined by specifying a nonstandard prolongation $\{G_n\}_{n<t}$ of the original wide sequence, non standard n < t and a certain Boolean valued extension $K(F)$ of $\cal M$. 

When the original sequence is definable in N (as it is always the case in this
thesis) then we can take for the prolongation the sequence defined
in the same way in M (and it is, in fact, also unbounded); hence
we will not refer explicitly to the choice of the prolongation.

On $K(F)$ we define a particular binary Boolean valued relation
$\Gamma$ (depending on n). Our aim is to relate first-order properties
of $\Gamma$ to various combinatorial and complexity-theoretic
properties of the wide sequence. In this case we shall talk about 
first order limit.

 We shall be also interested in second order properties of $\Gamma$,
for example the existence of a large clique. For this we expand
$K(F)$ to $K(F,G)$ with $G$ being second order objects (relations, functions).
In this case we shall talk about second order limit.

 

\section{The first order wide limit}

From now on we closely follow Chapter 1 of \cite{krajicek2010forcing}. Let $\mathcal{M}$ be the $\aleph_1$-saturated model of true arithmetic discussed in the previous chapter and let $\mathcal{G}_k$ be a wide sequence of graphs and $\Omega:=\mathcal{G}_n$ for $n\in\mathcal{M}\setminus \mathbb{N}$. One can check that graphs in $\Omega$ are all pseudofinite. The model $\M$ treats all its elements (including those which represent sets) as ``finite objects'' which lets us define uniform probability even on sets which are infinite from the set-theoretical perspective. We provide the definition generally for any $L$-reduct of $\M$, where $L\subseteq L_{all}$, but in the rest of the thesis we assume $L=L_{all}$ even though some finite sublanguage would suffice.


\begin{defi}
Let $\mathcal{A}:=\{A\in\mathcal{M};A\subseteq \Omega\}$ be the set of all subsets of $\Omega$ represented by an element in $\M$.

We define the \textbf{counting measure} as the uniform probability of $A$ when we sample $\Omega$ uniformly, so we have
\[A\in\A \to \abs{A}/\abs{\Omega},\]
the counting measure takes values in $\M$-rationals.
\end{defi}

One can check that $\A$ is a Boolean algebra, but not a $\sigma$-algebra as it is not closed under all countable unions. Indeed all singleton sets are part of $\A$ but the set of all elements with standardly many predecessors in $\Omega$ is not in $\A$.

\begin{defi}\label{defiinfinitesimal}
We call an $\M$-rational \textbf{infinitesimal} if it is smaller than all standard fractions $\frac{1}{k},k\in\N$.

Define an ideal in $\A$ as $\I:=\{A\in\A;\abs{A}/\abs{\Omega}\text{ is infinitesimal}\}$. Define the boolean algebra $\B:=\A/\I$. The induced measure on $\B$ is a real-valued measure and can be written as \[\mu(A/\I)=\st(\abs{A}/\abs{\Omega}).\]
\end{defi}

We can also check, that now $\mu$ is a measure in the ordinary sense and that $\B$ is an $\sigma$-algebra. In fact, the following key lemma holds.

\begin{lemm}
$\B$ is a complete boolean algebra.
\end{lemm}

We now define what we require of the family of functions $F$ we already mentioned.

\begin{defi}
Let $F$ be a non-empty set of $\M$-finite function which are elements in $\M$. We call it a \textbf{random vertex family} if it satisfies the following:
\begin{itemize}
\item The domain of any function $\alpha\in F$ is $\Omega$ and the range is $\set{g_n}$.
\end{itemize}
\end{defi}

Note that while every $\alpha\in F$ is represented by some element in $\M$, this need not be the case for the whole family $F$. From now on $\Gamma$ denotes a binary relational symbol which we understand as `limit edge relation'. For a $\{E\}$-formula $\varphi$ we denote its translation to a $\{\Gamma\}$-formula by rewriting every $E$ to $\Gamma$ as $\varphi^{\Gamma}$.

Now we can finally define the first order wide limit.

\begin{defi}
We define a $\B$-valued $\Gamma$-structure $\lim_{k\to n}^{F}G_k$, where $\Gamma$ is a binary relation symbol, with universe $F$ and $\Gamma$-sentences being evaluated by the following inductive conditions:
\begin{itemize}
\item $\bbl \alpha=\beta \bbr := \{\omega\in\Omega; \alpha(\omega)=\beta(\omega)\}/\I.$
\item $\bbl \Gamma(\alpha,\beta) \bbr := \{\omega\in\Omega; E_\omega(\alpha,\beta))\}/\I$.
\item $\bbl-\bbr$ commutes with $\land$,$\lor$,$\lnot$.
\item $\bbl(\exists x)A(x)\bbr:=\Lor_{\alpha\in F}\bbl A(\alpha)\bbr$.
\item $\bbl(\forall x)A(x)\bbr:=\Land_{\alpha\in F}\bbl A(\alpha)\bbr$.
\end{itemize}

By abuse of notation we will usually denote the limit $\lim_{F}\G_n$.
\end{defi}

We will also define the structure $K(\G_n,F)$ which is not important in the second order case, but makes the definition of its second order counterpart more manageable. 


\begin{defi}
$K(\G_n,F)$ will denote a $\B$-valued $L_{rel}\cup\{\Gamma\}$-structure defined as an $L_{rel}\cup\{\Gamma\}$-expansion of $\lim_F\G_n$. The boolean evaluations of atomic $L$-sentences are defined by the following inductive rules:

\begin{itemize}
\item $\bbl R(\alpha_1,\dots,\alpha_k) \bbr := \{\omega\in\Omega; R(\alpha_1,\dots,\alpha_k)\}/\I$ for any $k$-ary $R\in L_{rel}$.
\end{itemize}
\end{defi}


\section{The second order wide limit}

While we can find a truth value of a sentence in the language of graphs in the limit $\lim_{F}\G_n$, we will encounter situations where this is not sufficient to analyze the wide sequence $\{\G_k\}_{k>0}$. 

In Chapter \ref{chapdense} we will investigate how the existence of large cliques corresponds to the size of cliques in the limit graph. First we need some way to witness subsets of vertices -- this leads us to the second order wide limit. However, in the second order case the underlying arithmetic is much more important because we cannot just measure the set-theoretical cardinality of any such clique. For specific $n$ we could very well have $\card(\set{\floor{\log{n}}})=\card(\set{\floor{\frac{n}{2}}})$ but from the point of view of complexity theory, cliques of size $\floor{\log n}$ and $\floor{\frac{n}{2}}$ are dramatically different. In other words, our goal is also to have means to count the number elements of subsets or relations with values in (random variables in) $\M$.

When we say second order we mean two-sorted first order structures where one sort represents the usual elements (the `first order' sort) and the other represents functions on those elements (the `second order' sort). The `second order' here can also represent sets and relations by \{0,1\} valued functions.

\begin{defi}
We call a set of functions $G\subseteq \M$ an $F$-compatible \textbf{functional family} if every $\Theta\in G$ assigns to every $\omega\in\Omega$ a function $\Theta_\omega\in\M$ and after we define
\[\Theta(\alpha)(\omega):=
\begin{cases}\Theta_\omega(\alpha(\omega))&\alpha(\omega)\in\dom(\Theta_\omega)\\0&\text{otherwise,}
\end{cases}\]
we have that for every $\alpha \in F$ and $\Theta \in G$ we have $\Theta(\alpha)\in F$.

We call $\Theta\in G$ a (graph) $G$-\textbf{relation} if for every $\omega\in\Omega$ we have for some $k>0$ that $\dom \Theta_\omega \supseteq \set{g_n}^k$ and $\Theta_\omega:\dom \Theta_\omega \to \{0,1\}$.
\end{defi}

\begin{defi}
Let $F$ be a vertex family and $G$ be an $F$-compatible functional family. We define the structure $K(\G_n,F,G)$ as a two sorted $L$-structure with sorts $F$ and $G$ interpreting $L_{rel}$-sentences as $K(\G_n,F)$ and treating the sort $G$ as follows. First for equality we let
\[\bbl\Theta=\Xi \bbr := \{\omega\in \Omega;\Theta_\omega=\Xi_\omega\}/\I\]
and for the second order quantifiers we have the following inductive clauses
\begin{itemize}
\item $\bbl(\exists X)A(X)\bbr:=\Lor_{\Theta\in G}\bbl A(\Theta)\bbr$
\item $\bbl(\forall X)A(X)\bbr:=\Land_{\Theta\in G}\bbl A(\Theta)\bbr.$
\end{itemize}

We define \textbf{the second order wide limit} $\lim_{F,n}^G\{\G_k\}_{k>0}$ as $\{\Gamma\}$-reduct of $K(\G_n,F,G)$. By abuse of notation we will mostly denote it $\lim_F^G\G_n$.

(TODO: Terms)
\end{defi}

Let us note that if we have multiple boolean valued structures $\A_1,\A_2,\dots$ we may add the name of the structure as a prefix to the evaluation function to get $\A_1\bbl-\bbr$ and $\A_2\bbl-\bbr$ to avoid ambiguity or to emphasize the structure where the evaluation takes place. This is different from the standard notation, which would be including the structure name in the superscript as $\bbl\dots\bbr^{\A_1}$, but it is in our case preferable for typographical reasons.
 
\section{The vertex family $F_{rud}$ and $G_{rud}$}\label{secFrud}

Throughout this thesis we will mostly work with the vertex family $F_{rud}$ which ties the properties of $\lim_F \G_n$ with decision tree complexity -- but decision trees can also be seen as an abstraction of oracle machine time we mentioned in the Introduction.

After we choose the sequence $\{\G_k\}_{k>0}$ and $n>\N$ we again put $\Omega:=\G_n$ and define $F_{rud}$ as follows.

\begin{defi}
We define a \textbf{decision tree} to be a binary tree $T\in\M$ with a labeling of vertices and edges $\ell$. The non-leaf vertices are labeled by pairs of numbers $(u,v)$, where $u,v\in\set{g_n}$ and each edge is labeled either by $1$ or $0$. Each leaf vertex is then labeled by some element of $\M$.

Every sample $\omega\in\Omega$ uniquely determines a path in $(T,\ell)$ by interpreting the vertex labels as ``is $(u,v)\in E_\omega$?'' and the edge labels as true $(1)$ and false $(0)$ and the path then uniquely determines an output.

We define $F_{rud}$ to be the set of all functions computed by some $(T,\ell)$ of depth at most $g_n^{1/t}$.
\end{defi}

%Note that if we are given a graph $G\in G_k$ on $g_k$ vertices, we need $l:=2\log g_n$ bits to represent it. In the pseudofinite case, $k=n$, if we are restricted to inspect/query at most $g_n^{1/t}$ edges then it corresponds to $2^{O(l/t)}$ many queries or the subexponential oracle time.

One can verify that $F_{rud}$ is an vertex family for any $L\subseteq L_{all}$.

The definition of $G_{rud}$ is a bit more involved. The functionals in it will be computed by tuples of elements from $F_{rud}$ in the following sense.

\begin{defi}
Let $\hat\beta=(\beta_0,\dots,\beta_{m-1})\in\M$ be a $m$-tuple of elements in $F_{rud}$, for any $\alpha\in F_{rud}$ and $\omega\in\Omega$ we define
\[\hat\beta(\omega)=
\begin{cases}
\beta_{\alpha(\omega)}(\omega)&\alpha(\omega)<m\\
0&\text{otherwise.}
\end{cases}\]
\end{defi}

\begin{defi}
The family $G_{rud}$ consists of all functionals $\Theta$ such that there is $m\in\M$ and some $\hat \beta=(\beta_0,\dots,\beta_{m-1})\in\M$ that computes it.
\end{defi}

\begin{lemm}
$G_{rud}$ is $(F_{rud})$-compatible.
\end{lemm}
\begin{proof}
By induction in $\M$ we have that all the depth of all the trees is bounded by $g_n^{1/t}$ for some $t>\N$.

If we take some $\Theta\in G_{rud}$ and $\alpha\in F_{rud}$ we can compute $\Theta(\alpha)$ also by a tree in $F_{rud}$ by concatenating the trees computing $\alpha$ and $\beta_i$s.
\end{proof}

\section{Different choices of $n$}

Even though we generally pose no requirements on $n>\N$ there are examples of wide sequences for which the limit is sensitive on the choice of the non-standard number $n$.

\begin{exam}
Let 

\[\
G_{k}:=
\begin{cases}
\{(\set{k},E);\abs{E}=2,E(0,1)\}&\text{$k$ even}\\
\{(\set{k},E);\abs{E}=1,\lnot E(0,1)\}&\text{$k$ odd.}\\
\end{cases}\]

Let $n>\N$ then
\begin{align}
\lim_{F_{rud}}\G_{2n+1} \bbl \Gamma(0,1)\bbr =\0,
\end{align}
but
\begin{align}
\lim_{F_{rud}}\G_{2n}\bbl \Gamma(0,1)\bbr=\1.
\end{align}

Even though the concrete wide limits we will investigate in the following chapters do not depend on the specific $n$, it is important to note that we cannot generally remove the parameter $n$ from the definition of the limit.
\end{exam}

\section{Theories of wide limits}

If $\lim_F\G_n$ is the first order wide limit we will be interested in which exact $\{\Gamma\}$-sentences are valid in it. By \textbf{valid} we mean that their $\bbl\dots\bbr$ value is $\1$.

\begin{defi}
We define $\Th(\lim_F\G_n)$ as the set of all valid $\Gamma$-sentences in $\lim_{F}\G_n$.
\end{defi}

In the next chapter (Theorem \ref{thrmuni}) we will see that if a universal $\{E\}$-sentence $\varphi$ holds for all $G\in\G_k$ for $k$ big enough then $\lim_F\G_n\bbl\varphi^\Gamma\bbr=\1$. In particular, a wide limit of undirected graphs in a boolean valued undirected graph and a wide limit of directed graphs is a boolean valued directed graph.

Lastly, let us recall the concept of 0-1 laws which say that a certain probability tends either to $0$ or $1$ and not to any intermediate value. Here, instead of probability, we can think about the boolean values $\0$ and $\1$ and ask when does it happen that $\lim_F\G_n\bbl\varphi\bbr\in\{\0,\1\}$ for every $\{\Gamma\}$-sentence $\varphi$. This is exactly equivalent to the situation where the theory $\Th(\lim_G \G_n)$ is complete. Later in the thesis we prove such 0-1 law for several wide limits.

\chapter{The limit}
\section{Basic examples}
\subsection{$\mathcal{G}_k=\EDGE_k$}
We consider the classes of graphs
\[\EDGE_k:=\{\text{$(k,E); \abs{E}=1$}\},\]
and we let $\mathcal{G}_k=\EDGE_k$ and $F=F_{tree}$.

Intuitivelly one should not be able to find the edge on a significant i.e. non-infinitesimal fraction of samples with a tree that is allowed to explore only an infinitesimal fraction of edges.

\begin{thrm}\label{thrmedge}
We will prove that \[\bbl(\exists x)(\exists y)\Gamma(x,y)\bbr=\textbf{0}.\]
\end{thrm}
\begin{proof}
Let $T\in F_{tree}$ be a tree of depth $n^{1/t}$, for some $t>\mathbb{N}$ that outputs a pair of numbers less than $n$.

Start from the root of $T$ and always choose the path that corresponds to an edge not existing. At the end we obtain some answer, that gives us a set of at most $2 \cdot n^{1/t}+2$ vertices. Now we can find at least: 
\begin{align}
\binom{n-2n^{1/t}-2}{2}&=\frac{(n-2n^{1/t}-2)(n-2n^{1/t}-3)}{2}
\end{align}
different $\omega\in\Omega$ such that $T(\omega)$ is not an edge in $\omega$. 

The probability that any of those graphs is sampled is
\begin{align}
\frac{\binom{n-2n^{1/t}-2}{2}}{\abs{\mathcal{G}_k}}&=\frac{(n-2n^{1/t}-2)(n-2n^{1/t}-3)}{n(n-1)}\\
&=\left(1-\frac{2n^{1/t}-2}{n}\right)\cdot\left(1-\frac{2n^{1/t}-3}{n-1}\right)\\
&\geq \left(1-\frac{2n^{1/t}-2}{n}\right)^2\\
&\geq 1-\frac{4n^{1/t}-4}{n}.
\end{align}
And one can clearly see that $\st(1-\frac{4n^{1/t}-4}{n})=1$. This proves that the boolean value we are considering is $\textbf{0}$ since we can the two witnesses for $x$ and $y$ into a tree that could find an edge with depth $n^{1/t}$ for some $t>\mathbb{N}$.
\end{proof}

\subsection{Sparse $\mathcal{G}_k$}

One can see that in Theorem \ref{thrmedge} the exact shape of graphs in $\mathcal{G}_k$ does not play a crucial role. If $\mathcal{G}_k$ consisted of all graphs on $k$ verticies containing say exactly one triangle, or any other fixed subgraph of constant size, and no other edges, we would still find that the non-existence is valid in the limit graph.

A more general case would be to consider a family of graphs in which there is an infinitesimally small chance that two independent uniformly random verticies have an edge between. However, this is not sufficient.

\begin{exam}
Let $\mathcal{G}_k$ consist of those graphs on the vertex-set $k$ which contain the edge $E(0,1)$ and then has exactly one other edge. 

As $k$ increases, the number of edges get smaller than any standard positive fraction. But
\[\bbl(\exists x)(\exists y)E(x,y)\bbr=\1,\]
as witnessed by $x$ being the constant $0$ and $y$ the constant $1$ both of which are computable by a tree of depth $0$.
\end{exam}

One can see that having distinguished verticies can ruin the sparseness implying the non-existence of edges in the limit graph. We want to distiguish from this situation by considering the sequences $\mathcal{G}_k$ satisfying the following natural definition.

\begin{defi}
We say that $\{\mathcal{G}_k\}_{k=0}^\infty$ is \textbf{isomorphism closed}, if there is $k_0$ such that for every $k>k_0$ if we have that $G_1\in\mathcal{G}_k$, $V_{G_2}=g_k$ and $G_1\cong G_2$ then $G_2\in\mathcal{G}_k$.
\end{defi}


\begin{thrm}\label{thrmsparse}
Let an isomorphism closed $\mathcal{G}_k$ have the following property. There is a sequence $\{b_k\}_{k}$ and for big enough $k$, a uniformly sampled $2$-element $\{u,v\}\subseteq g_k$ and every $G\in\mathcal{G}_k$ we have
\[\Pr[E_G(u,v)]\leq b_k,\]
and some $k_0$ such that $\lim_{k\to\infty}k^{1/k_0}b_k = 0.$ Then

\[\bbl(\exists x)(\exists y)\Gamma(x,y)\bbr=\0.\]
\end{thrm}
\begin{proof}
Let us define the number $c_{u,v}:=\abs{\{G\in\mathcal{G}_k; E_G(u,v)\}}$, which is the number of graphs $G$ in $\mathcal{G}_k$ satisfying $E_G(u,v)$. Of course $c_{u,u}=0$ for every $u$.

\vspace{0.5em}
\textbf{Claim:} Let $u\neq v,u'\neq v'$ be vertices, then $c_{u,v}=c_{u',v'}$.\\
\textit{proof of claim:} Let $\rho:=(u\:u')(v\:v')$ be a permutation with cycles $(u\:u')$ and $(v\:v')$. We can let $\rho$ act on $\mathcal{G}_k$ by sending $G$ to a graph $\rho(G)$ which renames the edges coordinate-wise.

The fact that $\mathcal{G}_k$ is isomorphism closed implies that $\rho$ restricts to a bijection:
\[\rho':\{G\in\mathcal{G}_k;E_G(u,v), \lnot E_G(u',v')\}\to \{G\in\mathcal{G}_k;E_G(u',v'),\lnot E_G(u,v)\}\]
which proves the claim. \qed
\vspace{0.5em}

Now we define a matrix with entries
\[a_{G,\{u,v\}}:=\begin{cases}1&E_G(u,v)\\0&\text{otherwise}\end{cases}\]
where the rows are indexed by one of $\abs{\mathcal{G}_k}$ many graphs in $\mathcal{G}_k$ and the columns are indexed by the $\frac{k(k-1)}{2}$ many $2$-element sets of numbers in $k$. We take any distinct vertices $u,v$ and define $p:=\Pr_{G\in \mathcal{G}_k}[E_G(u,v)]=\frac{c_{u,v}}{\abs{\mathcal{G}_k}},$ by the claim the choice of $u,v$ does not matter.

The assumption from the statement is equivalent to the equality 
\[\sum_{\{u,v\}}a_{G,\{u,v\}}\leq \frac{k(k-1)}{2}b_k\] 
for every $G$. We combine this with the claim and the definition of $p$ to get
\begin{align}
\frac{k(k-1)}{2}\abs{\mathcal{G}_k}p&=\sum_{\{u,v\}}\sum_{G\in\mathcal{G}_k}a_{G,\{u,v\}}\\
&=\sum_{G\in\mathcal{G}_k}\sum_{\{u,v\}}a_{G,\{u,v\}}\\
&\leq \abs{\mathcal{G}_k}\frac{k(k-1)}{2}b_k
\end{align}
which implies
\[p\leq b_k.\]

Now let $k:=n$ and let $T\in F_{tree}$ be a tree of depth $n^{1/t}$ for some $t>\mathbb{N}$, where every leaf of $T$ is labeled by some edge. Walk down the tree $T$ from the root by answering negatively to every edge, this gives us a set $E_T$ of all edges $T$ inspected or outputed and $\abs{E_T}\leq n^{1/t}+1$.

Now we just need to prove that the probability $T$ find an edge is infinitesimally small. This is enough to prove the theorem, since the trees computing any two witnesses for $x$ and $y$ in the statement can be combined to construct $T$ and if any tree $T$ succeeds with only infinitesimally small probability, no random vertices can witness an edge on a set of non-zero measure.

We use the fact that $p\leq b_n$ to derive 
\begin{align}
\Pr_{G\in\mathcal{G}_n}[T\text{ finds an edge}]&\leq \sum_{\{u,v\}\in E_T} \Pr_{G\in \mathcal{G}_n}[E_G(u,v)]\\
&=\sum_{\{u,v\}\in E_T} \frac{c_{u,v}}{\abs{\mathcal{G}_n}}\\
&\leq\sum_{\{u,v\}\in E_T} p\\
&=(n^{1/t}+1)p\\
&\leq(n^{1/t}+1)b_k\\
&\leq n^{1/k_0}b_k\\
&\approx 0,
\end{align}
which proves the theorem.
\end{proof}

The assumption $\lim_{k\to \infty} k^{1/k_0}b_k=0$ for some $k_0$ may seem unintuitive at first. However, it precisely reflects what is ``sparse" for the trees in $T_{tree}$. The following example shows that without the assumption the theorem fails.

\begin{exam}
Let $\mathcal{G}_k$ consist of all graphs on the vertex set $\set{k}$ with precisely $\lceil\frac{k(k-1)}{2\log k}\rceil$ edges.

Then we claim that
\[\bbl(\exists x)(\exists y)\Gamma(x,y)\bbr=\1.\]

Let $\alpha$ and $\beta$ be verticies computed by the tree of the same shape which inspects a set of any $n^{1/t}$ distinct edges for some $t>\mathbb{N}$. If it finds an edge we define $\alpha$ and $\beta$ in any way so they are the distinct verticies incidental with this edge. Otherwise we let $\alpha(\omega)=\beta(\omega)=0$.

Let $T$ be a tree of the same shape, that computes the pair $\{\alpha,\beta\}$ then we can compute the probability where such a tree fails as the fraction of all graphs which have no edges that $T$ inspects. Let $m=\binom{n}{2}$. We get

\begin{align}
\Pr_{G\in\mathbb{G}_n}[T\text{ fails}]&=\frac{\binom{m-n^{1/t}}{\ceil{ \frac{n(n-1)}{2\log n}} }}{\binom{m}{\ceil{ \frac{n(n-1)}{2\log n}}}}\\
&=\frac{\frac{(m-n^{1/t})!}{\ceil{\frac{n(n-1)}{2\log n}} ! (m-\ceil{ \frac{n(n-1)}{2\log n}}-n^{1/t})!}}{\frac{m!}{\lceil \frac{n(n-1)}{2\log n}\rceil!(m-\ceil{\frac{n(n-1)}{2\log n}})!}}\\
&=\frac{(m-n^{1/t})!(m-\ceil{ \frac{n(n-1)}{2\log n}})!}{m!(m-\ceil{\frac{n(n-1)}{2\log n}}-n^{1/t})!}\\
&=\prod_{i=0}^{n^{1/t}-1}\frac{m-\ceil{\frac{n(n-1)}{2}}-i}{m-i}\\
&=\left(1-\frac{\ceil{\frac{n(n-1)}{2\log n}}}{\frac{n(n-1)}{2}}\right)^{n^{1/t}}\\
&\leq\left(1-\frac{\ceil{\frac{n(n-1)}{2\log n}}}{\frac{n(n-1)}{2}}\right)^{n^{1/t}}\\
&\leq\left(1-\frac{1}{\log n}\right)^{n^{1/t}}
\end{align}

And for any standard $k$ we have
\begin{align}
\left(1-\frac{1}{\log n}\right)^{n^{1/t}}&\leq\left(1-\frac{1}{\log n}\right)^{k\cdot \log n}\\
&\leq (e^{-\frac{1}{\ln 2}})^k.
\end{align}

So $\st(\Pr_{G\in\mathcal{G}_n}[T\text{ fails}])=0$ and we get 
\begin{align}
\mu(\bbl(\exists x)(\exists y)\Gamma(x,y)) \bbr &\geq \mu(\bbl \Gamma(\alpha,\beta)\bbr)\\
&=\st(1-\Pr_{G\in\mathbb{G}_n}[T\text{ fails}])\\
&=1.
\end{align}
\end{exam}

\section{Isomorphism closed $\mathcal{G}_k$}

So far the measure of every truth value we encountered was either $0$ or $1$. Is there a sequence $\mathcal{G}_k$ whose $F_{tree}$-limit and an $\{E\}$-sentence $\varphi$ such that ${0<\bbl\varphi^\Gamma\bbr<1?}$ As in the case of edge existance for a limit of sparse graphs, it is not hard to come up with an example if we allow distiguishing elements in $\mathcal{G}_k$.

\begin{exam}
Let \[\mathcal{G}_k=\{(\set{k},E);\text{$E$ has exactly two edges, one of them being $\{0,1\}$}\},\]

and let $F=F_{tree}$. Then $\mu(\bbl(\exists x)(\exists y)\Gamma(x,y)\bbr)=\frac{1}{2}$.
\end{exam}
\begin{proof}
Let $T_0$ be a tree that always outputs $0$ and $T_1$ be a tree that always outputs $1$. We can prove that $\bbl\Gamma(0,1)\bbr\geq\bbl\Gamma(\alpha,\beta)\bbr$ for any $\alpha,\beta$. 
\end{proof}

For the case of isomorphism closed $\mathcal{G}_k$ we prove that every $\{E\}$-sentence has truth value either $\0$ or $\1$. We start with the existential case.

\begin{exam}
NOT TRUE! Consider all graphs with one edge, or one non-edge. $\mu(\bbl(\exists x)(\exists y)\Gamma(x,y)\bbr)=\frac{1}{2}.$
\end{exam}

\begin{thrm}
Let $\varphi_0(\overline x)$ be an open $\{E\}$-formula, and let $\mathcal{G}_k$ be isomorphism closed and categorical then
\[\bbl(\exists \overline x)\varphi_0^\Gamma(\overline x)\bbr\in\{\0,\1\}.\]
\end{thrm}

\begin{proof}
Let $\overline T$ be a tuple of trees computing $\overline \alpha$ such that 
\[p:=\mu(\bbl\varphi_0(\overline \alpha)\bbr)>\0.\]

We want to iterate $\overline T$ to amplify the probability of success. Let 
\[A=\{\omega \in \Omega; \omega \models \varphi_0(\overline\alpha(\omega)) \}\]
and $\pi\in S_{\abs{\G_n}}$ with maximal $\abs{\pi[A]\setminus A}$. Let $\pi(\overline T)$ be a tuple of relabeled trees computing $\pi(\overline \alpha)$.
\end{proof}

\chapter{Dense case}\label{chapdense}

\section{$\G_k=\SK_k^{1/2}$}

Now we turn to analyze dense wide sequences in the second order case; that is the case when we shall be interested in second order properties of the wide limit. In this chapter we assume $F=F_{rud}$ and $G=G_{rud}$. Specifically we will consider the problem of finding a large clique in a graph. 

Generally it is considered a computationally hard problem to find a large clique in a graph. It is an $\NP$-hard problem and thus it is conjectured that it cannot be solved in polynomial time. We first turn to the following wide sequence.

\begin{defi}
Let \[\SK_k^{1/2}=\{(\set{k},E);\text{$E$ consists of exactly one $k/2$-clique}\}.\]
\end{defi}

Limiting inputs to $\SK_k^{1/2}$ makes the problem less complex, because for a vertex $v$ to be a part of the biggest clique it is enough that it has nonzero degree. Naturally we want to see whether there is a large clique in $\lim_F^G \SK_n^{1/2}$, every sample has a clique of size $n/2$, but is there a set of $F$-vertices witnessing that? 

Here to measure the size of such a clique, the second order wide limit  by itself is not a sufficient object. Instead we need to turn to the arithmetical expansion $K(\G_n,F,G)$ and ask whether we can find an injective function from some large initial segment into a clique. This can be expressed as a second order sentence with parameters from $F$. It is not hard to prove the following result which also demonstrates how the sentence is formed.

\begin{thrm}\label{thrmsk1}For every $t>\N$ we have
\[K(\SK_k^{1/2},F,G)\bbl\text{$E$ has a clique of size $n^{1/t}$}\bbr=\1.\]

\end{thrm}

\begin{proof}(sketch) We need to analyze the value
\begin{align}
&\bbl(\exists \Lambda)(\forall u)(\forall v)(((u\not = v \leq n^{1/t})\to E(\Lambda(u),\Lambda(v)))\land(\Lambda:[n^{1/t}]\hookrightarrow\mathcal{M}))\bbr
\end{align}
which is equal to
\begin{align}
\Lor_\Lambda\Land_u\Land_v\bbl(u\not=v < n^{1/t})\to(E(\Lambda(u),\Lambda(v))\land\Lambda(u)\not=\Lambda(v))\bbr.
\end{align}
So we want to find some $n^{1/t}$-tuple of trees computing some $\Lambda$, computed by some $(\gamma_0,\dots,\gamma_{n^{1/t}-1})$, which is injective on $\set{n^{1/t}}$ and its $\set{n^{1/t}}$ range is a clique in $E$.

We define $(\gamma_0,\dots,\gamma_{n^{1/t}-1})$ as follows. The tree $T_0$ computing $\Lambda_0$ inspects all the edges $(u,v)\in\set{n^{1/t}}\times\set{n^{1/t}}$ in some specified order and outputs the first vertex it finds with an edge and outputs $0$ otherwise. The tree $T_i$ computing $\gamma_i$ extends $T_{i-1}$ by not outputting the $i$-th vertex with the property but instead continues inspecting and outputs the $(i+1)$-th vertex for which it finds an edge. Also every such tree has depth at most $n^{1/t}\cdot n^{1/t}=n^{2/t}=n^{1/(t/2)}$.

One can check that the probability the tree $\gamma_i$ does not find $i$ vertices with an edge is infinitesimal and therefore it always outputs a vertex in the clique of $\omega$. Intuitively this is because the expected ratio of the vertices in $\set{n^{1/t}}$ which are also in the clique is about one half. Moreover, every $\gamma_i$ outputs the $i$-th element element of the ordering and thus it is injective.
\end{proof}

At first glance the lower bound $n^{1/t}$ for every nonstandard $t$ may seem optimal given the proof method we used, but there is a way to radically improve it. The idea is to partition the set of vertices into many smaller ones and let $\gamma_i$ search only in the $i$-th set. First we need the following lemmas.

\begin{lemm}\label{lemmmiss}
Let $S\subseteq \set{n}$ such that $\abs{S}=m > \N$, then
\[\st\left(\Pr_{G\in\G_n}[\text{$S$ contains no vertices in the clique of $G$}]\right)=0.\]
\end{lemm}
\begin{proof}
There are $\binom{n-m}{\floor{\frac{n}{2}}}$ different graphs in $\SK_n^{1/2}$ in which the clique does not intersect $S$. We then bound the probability as
\begin{align}
\frac{\binom{n-m}{\floor{\frac{n}{2}}}}{\binom{n}{\floor{\frac{n}{2}}}}&=\frac{(n-m)!(n-\floor{\frac{n}{2}})!}{(n)!(n-\floor{\frac{n}{2}}-c)!}\\
&=\prod_{i=0}^{m-1}\frac{(n-i-\floor{\frac{n}{2}})}{(n-i)}\\
&=\prod_{i=0}^{m-1}\left(1-\frac{\floor{\frac{n}{2}}}{n-i}\right)\\
&\leq \left(1-\frac{\floor{\frac{n}{2}}}{n}\right)^m\\
&\leq \left(1-\frac{\floor{\frac{n}{2}}}{n}\right)^{n \cdot \frac{m}{n}}\\
&\leq e^{-\floor{\frac{n}{2}}\frac{m}{n}}.
\end{align}
But $\floor{\frac{n}{2}} \frac{m}{n}$ is infinite therefore the bound is infinitesimal.
\end{proof}

\begin{lemm}\label{lemmfindtree}
Let $a\in\M$, let $v_0,\dots,v_{a-1}\in \set{n}$ distinct vertices, then there exist trees $T_{v_i}$ such that
\[\st\left(\Pr_{\omega\in\Omega}[\forall i: (v_i,T_{v_i}(\omega))\in E_\omega|\forall i: \text{$v_i$ is in the clique}]\right)=1.\]
\end{lemm}

\begin{proof}
The tree $T_{v_i}$ inspects all the edges $(v_i,j)$ where $j$ ranges over $\set{n^{1/t}}$ for some $t>\N$ and outputs $j$ if $(v_i,j)\in E_\omega$. By Lemma \ref{lemmmiss} we have that only an infinitesimal number of graphs have their clique not intersect $\set{n^{1/t}}$ so each $T_{v_i}$ succeeds on all but an infinitesimally small portion of $\Omega$. But if one $T_{v_i}$ finds a neighbour of $v_i$, then all do, since nonzero degree vertices in every $\omega$ form a clique and the same $w\in\set{n^{1/t}}$ is a neighbour of all $v_i$s.
\end{proof}

\begin{lemm}\label{lemmmiss2}
Let $S_0,\dots,S_{a-1}\subseteq\set{n}$ sets of size $a\in \M$ for $i\in \M$ then
\[\Pr_{\omega\in\Omega}\left[\bigcup_{i=0}^{a-1}[\text{$S_i$ contains no vertices in the clique of $\omega$}]\right]\leq a\cdot e^{-\floor{\frac{n}{2}}\frac{m}{n}}.\]
\end{lemm}
\begin{proof}Follows from the proof of Lemma \ref{lemmmiss} and union bound.
\end{proof}

Now we are ready to improve on Theorem \ref{thrmsk1}.

\begin{thrm}\label{thrmsk}
Let $m = 2 \ln n$. Then
\[K(\SK_k^{1/2},F,G)\bbl\text{$E$ has a clique of size $\floor{n/m}$}\bbr=\1.\]
\end{thrm}
\begin{proof}
Partition a subset of $\set{n}$ to sets $S_{0},\dots,S_{\floor{n/m}-1}$ each of size at least $m$. Specifically if $m$ divides $n$ then we partition the whole $\set{n}$.

By Lemma \ref{lemmmiss2} we have that with probability that we do not sample $\omega$ which have the clique intersect all $S_i$s 
\begin{align}
\floor{\frac{n}{m}}\cdot e^{-\floor{\frac{n}{2}}\frac{m}{n}} = e^{\ln\floor{\frac{n}{m}}-\floor{\frac{n}{2}}\frac{m}{n},}\label{eqpr}
\end{align}
we can bound the exponent as
\begin{align}
\ln\floor{\frac{n}{m}}-\floor{\frac{n}{2}}\frac{m}{n}&\leq\ln\frac{n}{m}-\frac{n}{2}\cdot\frac{m}{n}+\frac{m}{n}\\
&\leq\ln n - \ln m - \frac{m}{2} + \frac{m}{n}1\\
&\leq \ln n - \ln \ln n - \ln n+\frac{2\ln n}{n}\\
&\leq - \ln \ln n + \frac{2\ln n}{n}
\end{align}
which is negative and infinite, because $\frac{\ln x}{x}\overset{x\to \infty}{\to} 0$ as  and therefore \eqref{eqpr} is infinitesimal.

So with probability infinitesimally close to 1 we have in each $S_i$ a vertex $v_i$ which is also a part of the clique. By Lemma \ref{lemmfindtree} we have that there exists a tree verifying whether a given vertex is in the clique and since $m\leq n^{1/t}$ for some $t$ we can concatenate the trees to get a tree $T_{S_i}$ which finds in $S_i$ an element of the clique with probability infinitesimally close to $1$.

Finally we can have a function $\Lambda\in G$ computed by $(\Lambda_0,\dots,\Lambda_{\floor{n/m}-1})$ by letting $\Lambda_i$ be computed by $T_{S_i}$ we have already verified $\bbl\Lambda\text{ is a clique}\bbr=\1.$

Because $T_{S_i}(\omega)\in S_i$ when $S_i$ succeeds, and $S_i$ are disjoint we have \[\bbl\Lambda:\set{\floor{n/m}}\xhookrightarrow{} \M\bbr=\1. \]
Which proves the theorem.
\end{proof}

One can also check that for $\SK_{k}^{1/l}$, the graphs whose edges are exactly one $\floor{k/l}$ clique for some $l\in \N$, the wide limit has a clique of size $\floor{n/m}$ for $m=l\cdot \ln(m)$ by the same technique.

Even though the size of the clique has radically increased, we still did not find a clique in $E$ of size $n/2$. One can verify that with the method provided one cannot obtain such a clique because the probability that any of $\floor{n/2}$ two-element sets does not intersect the clique is too large. 

\begin{ques}
Is $\floor{n/(2\ln n)}$ the clique number in the Boolean valued graph $\lim_{F_{rud}}^{G_{rud}}\SK_n^{1/2}$, measured in $K(\SK_k^{1/2},F,G)$?
\end{ques}


\section{$\G_k=\CK_k^{1/2}$}

Now let us mention the more complex case of the wide sequence $\CK_k^{1/2}$
\[\CK_k^{1/2}=\{(\set{k},E);E\text{ contains a $\floor{k/2}$ clique.}\}\]

We are still guaranteed that every $\omega$ contains a large clique, but there is no easy way to check whether a given vertex $v$ is contained in the large clique. To prove the following theorem we can use that $\CK_k^{1/2}$ is isomorphism closed to translate the case for $\SK_k^{1/2}$ at least for the case of cliques of standard size.

\begin{thrm}\label{thrmck}
Let $c\in\N$ then 
\[K(\CK_k^{1/2},F,G)\bbl\text{$E$ has a clique of size $c$}\bbr>\0.\]
\end{thrm}
\begin{proof}
Consider the set \[\G_n^+=\{(G,U);\text{$G\in \G_n$ and $U\subseteq \set{n}$ is a $\floor{n/2}$ clique in $G$ }\}\]
and projections
\begin{align*}
\pi_1:\G_n^+&\to\G_n,(G,U)\mapsto G.\\
\pi_2:\G_n^+&\to\Po(\set{n}),(G,U)\mapsto U
\end{align*}

From the fact that $\CK_k^{1/2}$ is isomorphism closed, we have that $\abs{\pi_2^{-1}[U_1]}=\abs{\pi_2^{-1}[U_2]}$. We will set $\nu:=\abs{\pi_2^{-1}[U]}$ for some clique $U$ of size $\floor{n/2}$ so we have $\nu\abs{\G_n}=\G_n^+$. Let $\varphi_c(\overline x)$ be a $\{E\}$-sentence is satisfied iff $\overline x$ is a $c$ sized clique.  Denote an event in the sample space $\G_n$
\[A_{\overline a} :=\{G\in\G_n;G\models \varphi_c(\overline a)\}.\]
and in the sample space $\G_n^+$
\[A_{\overline a}^+ :=\{(G,U)\in\G_n^+;\overline a\text{ is a subclique of $U$}\}\]
notice that if $(G,U)\in A_{\overline a}^+$ then $G\in A_{\overline a}$ which implies (by another argument with projections) that
\[\nu\cdot\abs{A_{\overline a}}\geq \abs{A_{\overline a}^+}\]
so in conjunction with Example \ref{examskhalf} we have
\begin{align*}
\Pr_{G\in\G_n}[A_{\overline a}]&\geq \frac{\abs{A_{\overline a}^+}}{\abs{\G_n^+}}\\
&= \Pr_{G\in\SK_{k}^{1/2}}[\overline a \text{ is in the clique}]\\
&>(1/2)^c
\end{align*}
and since $c$ is constant we have $\Pr_{G\in\G_n}[A_{\overline a}]>0$.
\end{proof}

Of course it remains to show that using trees can actually increase the value all the way to $\1$, we did not get to prove that.

Finally one can return to cliques of nonstandard size and intuitively one expects this to be hard. For a clique of size greater than $n^{1/t}$ for any $t>\N$ one has to check more than $n^{1/t}$ edges to even know whether a given set is a clique and the counting argument used in the previous theorem implies that just guessing the clique of nonstandard size is not enough. We therefore present the following conjecture.

\begin{conj}
Let $m\leq c \ln n$ for any $c\in\N$, then
\[K(\CK_k^{1/2},F,G)\bbl\text{$E$ has a clique of size $\floor{n/m}$}\bbr=\0.\]
\end{conj}

\chapter{Sparse case and \texorpdfstring{$\TFNP$}{TFNP}}\label{chapsparse}

\section{\texorpdfstring{$\mathcal{G}_k=*\text{PATH}_k$}{Gk=*PATHk}}

Now we turn our attention to a wide sequence which is made up of `the hardest instances in $\LEAF$'. That is, if we are given a degree 1 vertex labeled $0$ and search for another degree 1 vertex, it is the hardest if there are only two degree 1 vertices and the path from $0$ to the solution is as long as possible.

\begin{defi}
We define $\pPATH_k$ (the pointed paths on $k$ vertices) as the set of all (undirected) graphs $G$ on the vertex set $\set{k}$, where $G$ is isomorphic to the path on $k$ vertices and $\deg_G(0)=1$.
\end{defi}

One can also see that it is not fruitful to analyze the $F_{rud}$ limit, because there are only $k-1$ edges in each $G\in\pPATH_k$ and therefore, by an analogous argument to the proof of Theorem \ref{thrmsparse}, we have that $\lim_{F_{rud}}\pPATH_k\bbl(\exists x)(\exists y)E(x,y)\bbr=\0.$ Moreover in the type 2 version of the problem the graph is presented by an oracle which gives us the neighbour set for each vertex, so we define a new family $F$ as follows.

\begin{defi}
After we fix $n$, we define $F_{nbtree}$ as the set of all functions computed by some some labeled tree with the following shape:

\begin{itemize}
\item Each non-leaf node is labeled by some $w\in\set{n}$. 
\item For each $\{u,v\} \subseteq \set{n}$ and a node $w$ there is an outgoing edge from $w$ labeled $\{u,v\}$.
\item Each leaf is labeled by some $m\in \set{g_n}$.
\item The depth of the tree is at most $g_n^{1/t}$ for some $t>\mathbb{N}$.
\end{itemize}

Computation of such a tree on an undirected graph $G$ goes as follows. We interpret the non-leaf nodes as questions "what is the neighbour set of $w$?" and the edges as answers from our graph $\omega$, and thus we follow a path in the computation tree determined by $\omega$ until we find a leaf, in which case the computation returns the label of the leaf. We will denote the set of all such trees as $\T_{nbtree}$.
\end{defi}

We now shift our focus to analyzing the ability of trees from $F_{nbtree}$ to find the non-zero degree 1 vertex in $G\in \pPATH_n$. We say that a tree $T\in \T_{nbtree}$ fails at a graph $G$ if $T(G)$ is not a non-zero vertex of degree one in $G$.

\begin{defi}
Let $m\leq n$ and $v\in \set{w}$ and $U\subseteq \set{w}$ with $\abs{U}\leq 2$, then we define 

\[\mathcal{G}_m^{v?=U}:=\{G\in\mathcal{G}_m;N_G(v)=U\},\]
where $N_G$ is the neighbour-set function of $G$.
\end{defi}
\begin{lemm}
There are bijections for all nonstandard $m\leq n$ and distinct $u,v,w\in \set{m}\setminus\{0\}$:

\begin{align}
\mathcal{G}_m^{v?=\{u,w\}}&\cong\mathcal{G}_{m-2}\times [2] \label{pPathfirstbij}\\
\mathcal{G}_m^{v?=\{u,0\}}&\cong\mathcal{G}_{m-2}\label{pPathsecbij}\\
\mathcal{G}_m^{0?=\{u\}}&\cong\mathcal{G}_{m-1}\label{pPaththrdbij}.
\end{align}
\end{lemm}
\begin{proof} (sketch)
For \eqref{pPathfirstbij} we can just contract all of {u,v,w} into one vertex and relabel the rest of the graph, leaving the orientation as one remaining bit of information. This is obviously reversible and a bijection.

For \eqref{pPathsecbij} we can do the same, but the orientation is given by $0$.
\end{proof}

\begin{lemm}\label{lemmrelabeltree}
Let $T\in \T_{nbtree}$, with root labeled $v\in[m]\setminus{0}$, we have for each $T_{v?=\{u,w\}}$ a tree $\tilde T_{v?=\{u,w\}}$ of the same depth, such that
\begin{align}
\Pr_{G\in\G_m}(T_{v?=\{u,w\}}\text{ fails}|v?=\{u,w\})=\Pr_{G\in\G_{m-2}}(\tilde T_{v?=\{u,w\}}).
\end{align}
For a tree $T$ with the root labeled $0$, we have a tree $\tilde T_{v?=\{u,w\}}$ of the same depth, such that
\begin{align}
\Pr_{\G\in\G_m}(T_{v?=\{u\}}\text{ fails}|v?=\{u\})=\Pr_{\G_{m-1}}(\tilde T_{v?=\{u\}}).
\end{align}
\end{lemm}
\begin{proof}(sketch) To construct the tree, we just replace all vertices in labels of $T_{v?=\{u,w\}}$ by their renumbering from the bijection in \eqref{pPathfirstbij}.

One can then check that the trees $T_{v?=\{u,w\}}$ and $\tilde T_{v?=\{u,w\}}$ are isomorphic in a sense that their computation of a graph $G$ and $\tilde G$ respectively, $\tilde G$ being the corresponding $(m-2)$-vertex graph, agree with the structure of the path and that correctness of leaves is preserved under the renumbering. Essentially they emulate the same computation but on a smaller graph.
\end{proof}


\begin{lemm}\label{lemmpPATHlowerbound}
For all nonstandard $t>\mathbb{N},m\geq n-2n^{1/t}$ and $k\in [n^{1/t}+1]$ for all trees $\T\in \T_{nbtree}$ of depth $k$ we have
\[
\Pr_{G\in\G_m}(\text{$T$ fails})\geq \prod_{i=0}^k\left (1-\frac{2}{m-2i-2}\right ).
\]
\end{lemm}
\begin{proof}
We proceed by induction on $k$. 

$k=0:$ We have that the probability of success of a straight guess is at most $\frac{1}{m-1}$. Therefore

\begin{align}
\Pr_{G\in\G_m}(\text{$T$ fails})\geq \left (1-\frac{1}{m-1}\right ) \geq \left(1-\frac{2}{m-2}\right).
\end{align}

$(k-1) \Rightarrow k:$ First we assume that the root is labeled $0$. Then we have 

\begin{align}
\Pr_{G\in\G_m}[\text{$T$ fails}]&=\sum_{u\in V\setminus\{0\}}\Pr_{G\in\G_{m-1}}[E(0,u)]\Pr_{G\in\G_{m-1}}[T_{0?=\{u\}}\text{ fails}|E(0,u)]\\
&\geq \Pr_{G\in\G_{m-1}}[T_{0?=\{u\}}\text{ fails}|E(0,u)]\\
&= \Pr_{G\in\G_{m-1}}[\tilde T_{0?=\{u\}}\text{ fails}]\\
&\geq \prod_{i=0}^{k-1}\left(1-\frac{2}{m-2i-2}\right)\\
&\geq \prod_{i=0}^{k}\left(1-\frac{2}{m-2i-2}\right).
\end{align}

Now we assume that the root is labeled $v\not= 0$. First we notice that

\begin{align}
\Pr_{G\in\G_m}[E(v,0)]&=\frac{1}{m-1}\\
\Pr_{G\in\G_m}[N(V)=1]&=\frac{1}{m-1}\\
\Pr_{G\in\G_m}[\abs{N(V)\setminus\{0\}}=2]&=1-\frac{2}{m-1},\label{pmmiddle}
\end{align}
the first two probabilities are obviously $\frac{1}{m-1}$ because they correspond to $v$ being positioned on one of the ends of the non-zero segment which has length $m-1$. The event in \eqref{pmmiddle} is the complement of the union of the first two events, which have empty intersection, giving us that stated probability.

Then we have for $p:=\Pr_{G\in\G_m}[\text{$T$ fails}]$
\begin{align}
p&=\Pr_{G\in\G_m}[E(v,0)]\Pr_{G\in\G_m}[\text{$T$ fails}|E(v,0)]\\
&\:\phantom{=}\:+\Pr_{G\in\G_m}[\abs{N(v)\setminus\{0\}}=2]\Pr_{G\in\G_m}[\text{$T$ fails}|\abs{N(v)\setminus\{0\}}=2]\\
&\:\phantom{=}\:+(\Pr_{G\in\G_m}[\abs{N(v)}=1]\Pr_{G\in\G_m}[E(v,0)]\Pr_{G\in\G_m}[\text{$T$ fails}|\abs{N(v)}=1])\\
&\geq \Pr_{G\in\G_m}[\abs{N(v)\setminus\{0\}}=2]\Pr_{G\in\G_m}[\text{$T$ fails}|\abs{N(v)\setminus\{0\}}=2]\\
&\geq (1-\frac{2}{m-1})\\
&\:\phantom{=}\:\sum_{\substack{u,w\in V\setminus \{0\} \\ u\not= w}}\Pr_{G\in\G_m}[v?=\{u,w\}]\Pr_{G\in\G_m}[T_{v?=\{u,w\}}\text{ fails}|v?=\{u,w\}]\\
&\geq(1-\frac{2}{m-1})\Pr_{G\in\G_m}[T_{v?=\{u_0,w_0\}}\text{ fails}|v?=\{u_0,w_0\}]\label{pmfailchoice}\\
&\geq(1-\frac{2}{m-1})\Pr_{G\in\G_{m-2}}[\tilde T_{v?=\{u_0,w_0\}}\text{ fails}]\label{pmfaillemma}\\
&\geq(1-\frac{2}{m-1})\prod_{i=0}^{k-1}(1-\frac{2}{m-2i-4})\label{pmfailih}\\
&\geq(1-\frac{2}{m-2})\prod_{i=1}^{k}(1-\frac{2}{m-2i-2})\\
&\geq\prod_{i=0}^{k}(1-\frac{2}{m-2i-2}).
\end{align}
where in \eqref{pmfailchoice} we choose $u_0,w_0$ with the lowest value of \[\Pr_{G\in\G_m}(\text{$T_{v?=\{u_0,w_0\}}|v?=\{u_0,w_0\}$}),\] the bound follows the fact that all $\Pr_{G\in\G_m}(v?=\{u,w\})$ are the same for distinct non-zero $u,w$. In \eqref{pmfaillemma} we use the lemma \ref{lemmrelabeltree} and in \eqref{pmfailih} we use the induction hypothesis.
\end{proof}

\begin{crll}\label{crllpPathtreefail}
For a tree $T\in \T_{nbtree}$ and $c\in \N$ we have that
\[\st\left(\Pr_{G\in\G_{n-c}}[T\text{ fails}]\right)= 1.\]
\end{crll}
\begin{proof}
Since $T$ has depth at most $n^{1/t}$ for some $t>\mathbb{N}$ we by the previous lemma that
\begin{align}
\Pr_{G\in\G_n}(\text{$T$ fails})&\geq \prod_{i=0}^{n^{1/t}}\left(1-\frac{2}{n-2i-c-2}\right)\\
&\geq \left(1-\frac{2n^{1/t}}{n-2n^{1/t}-c-2}\right)
\end{align}
and the standard part of the lower bound is $1$.
\end{proof}

Finally, we can prove the following theorem.

\begin{thrm}\label{thrmnoend}
\[\lim_{F_{nbtree}}\pPATH_n\bbl(\exists v)(\exists u)(\forall w)((v\not=0)\land (E(v,u)) \land(E(v,w) \to u=w))\bbr=\0\]
\end{thrm}
\begin{proof}
Expanding the value of the formula in the statement we get
\[\Lor_{\alpha}\Lor_{\beta}\Land_{\gamma}\bbl(\alpha\not=0)\land (E(\alpha,\beta)) \land(E(\alpha,\gamma) \to \beta=\gamma)\bbr,\]
to prove it evaluates to $\0$ we need to find for every $\alpha,\beta$ some $\gamma$ such that
\[\bbl(\alpha\not=0)\land (E(\alpha,\beta)) \land(E(\alpha,\gamma) \to \beta=\gamma)\bbr=\0.\]
For any $\alpha,\beta$ we define
\[\gamma(\omega):=\begin{cases}v&N(\alpha(\omega))=\{v\}\\w&w\in N(\alpha(\omega))\setminus\{\beta(\omega)\},\end{cases}\]
such a function can be computed by a tree in $F_{nbtree}$ which we can construct by concatenation of trees computing $\alpha$ and $\beta$.

Let $T$ be the tree computing $\alpha$. Now we proceed by contradiction, let
\[\epsilon:=\mu(\bbl(\alpha\not=0)\land (E(\alpha,\beta)) \land(E(\alpha,\gamma) \to \beta=\gamma)\bbr)>0,\]
by definition this means that
\[\epsilon=\st(P_n[(\alpha\not=0)\land (E(\alpha,\beta)) \land(E(\alpha,\gamma) \to \beta=\gamma)])>0.\]
Expanding the value of the formula in the statement we get
\[\Lor_{\alpha}\Lor_{\beta}\Land_{\gamma}\bbl(\alpha\not=0)\land (E(\alpha,\beta)) \land(E(\alpha,\gamma) \to \beta=\gamma)\bbr,\]
to prove it evalues to $\0$ we need to find for every $\alpha,\beta$ some $\gamma$ such that
\[\bbl(\alpha\not=0)\land (E(\alpha,\beta)) \land(E(\alpha,\gamma) \to \beta=\gamma)\bbr=\0.\]
For any $\alpha,\beta$ we define
\[\gamma(\omega):=\begin{cases}v&N(\alpha(\omega))=\{v\}\\w&w\in N(\alpha(\omega))\setminus\{\beta(\omega)\},\end{cases}\]
such a function can be computed by a tree in $F_{nbtree}$ which we can construct by concatenation of trees computing $\alpha$ and $\beta$.

Let $T$ be the tree computing $\alpha$. Now we proceed by contradiction, let
\[\epsilon:=\mu(\bbl(\alpha\not=0)\land (E(\alpha,\beta)) \land(E(\alpha,\gamma) \to \beta=\gamma)\bbr)>0,\]
by definition this means that
\[\epsilon=\st(P_n[(\alpha\not=0)\land (E(\alpha,\beta)) \land(E(\alpha,\gamma) \to \beta=\gamma)])>0.\]
But by definition of $\gamma$ and Corollary \ref{crllpPathtreefail} we have
\begin{align*}
0&<\epsilon\\
&=\st(P_n[(\alpha\not=0)\land (E(\alpha,\beta)) \land(E(\alpha,\gamma) \to \beta=\gamma)])\\
&\leq\st(P_n[(\alpha\not=0)\land (E(\alpha,\beta)) \land \abs{N(\alpha)}=1])\\
&\leq\st(P_n[(\alpha\not=0)\land \abs{N(\alpha)}=1])\\
&=\st(P_n[T\text{ does not fail}])\\
&=0.
\end{align*}
A contradiction.
\end{proof}

\begin{crll}\label{crllpPATH}
$\Th(\lim_{F_{nbtree}}\pPATH_n)$ is complete.
\end{crll}
\begin{proof}
By applying Theorem \ref{thrmuni} we have that the sentences $\lnot C_k$ stating that there are not cycles of length $k\in\N$ are valid all in $\lim_{F_{nbtree}}\pPATH_n$, the sentence $D^{1,2}_{1,rest}$ stating that there is exactly one vertex of degree $1$ and all other vertices have degree $2$ is valid by Theorem \ref{thrmnoend}. 

Let $T=\{\lnot C_k,k\in\N\}\cup\{D_{1,rest}^{1,2}\}$, and let $\A_1,\A_2\models T$, then we can see by the handshaking lemma that $\M_1$ and $\M_2$ are both infinite. And we can see that they can be decomposed into one path starting at $0$ with no end, and then more infinite paths which have the order type of $\Z$. The duplicator of Ehrenfeucht-Fraïssé game has a winning strategy by responding to any element on the order type $\Z$ with a far enough element on the path of the order $\N$.
\end{proof}

\section{\texorpdfstring{$\G_k=\pPATHl_k$}{Gk=*PATHlk}}

So far we have proved that the hardest instances of $\LEAF$ do not have a solution in the $F_{nbtree}$ limit and that they satisfy the 0-1 law. We can generalize this result to a larger class of instances.

\begin{defi}
We define $\pPATHl_k$ (the pointed paths on $k$ vertices of length at most $k$) as the set of all (undirected) graphs $G$ on the vertex set $\set{k}$, where $G$ has a subgraph isomorphic to the path on $l\leq k$ vertices, $\deg_G(0)=1$ and no other edges.
\end{defi}

Immediately we have that $\pPATH_k$ is a portion of $\pPATHl_k,$ we can prove even more.

\begin{defi}
We define $\pPATH_k^{\l}$ as the portion of $\pPATHl_k$ where the path subgraph is of length exactly $l$. 
\end{defi}


\begin{lemm}
Let $c\in \N$, then 
\[\lim_{k\to\infty}\frac{\abs{\pPATH_k^{k-c}}}{\abs{\pPATHl_k}}=\frac{1}{ec!}\]
So $\pPATH_k^{k-c}$ is a large portion of $\pPATHl_k$ and specially $\pPATH_k\leqp \pPATHl_k$.
\end{lemm}
\begin{proof}
By direct computation we have that the fraction $\abs{\pPATH_k}/\abs{\pPATHl_k}$ is
\begin{align}
\frac{(k-1)!/(c!)}{\sum_{i=1}^{k-1}\prod_{j=0}^{i-1}(k-j-1)}&=\frac{(k-1!)}{c!\sum_{i=1}^{k-1}\frac{(k-1)!}{(k-i-1)!}}\\
&=\frac{(k-1!)}{c!\sum_{i=1}^{k-c-1}\frac{(k-1)!}{(k-i-1)!}}\\
&=\frac{1}{c!\sum_{i=1}^{k-1}\frac{1}{(k-i-1)!}}\\
&=\frac{1}{c!\sum_{i=0}^{k-2}\frac{1}{i!}},
\end{align}
and the denominator tends to $ec!$ as $k\to\infty$.
\end{proof}

\begin{lemm}\label{lemmpathnnc}
Let $T\in \T_{nbtree}$ be a tree, then
\[\st\left(\Pr_{G\in\pPATH_{n}^{n-c}}[\text{$T$ fails}]\right )=1.\]
\end{lemm}
\begin{proof}(sketch)
By an analogous argument to how we proved Lemma \ref{lemmpPATHlowerbound} we get that
\[\Pr_{G\in\pPATH_{n}^{n-c}}[T\text{ fails}]\geq \prod_{i=0}^{n^{1/t}}\left(1-\frac{(c+2)}{n-2i-2}\right),\]
standard part of which is also $1$. The constant $c$ appears because there is a $(1-\frac{c}{m})$ chance of finding a degree $0$ in the induction step.
\end{proof}

\begin{lemm}\label{lemmpathmaj}
Let $\G_k':=\bigcup_{c\in\N}\pPATH^{k-c}_k$, then
\[\lim_{k\to\infty}\frac{\abs{\G_k'}}{\abs{\pPATHl_k}}=1.\]
In other words, $\G_k'$ is a major portion of $\pPATHl_k$.
\end{lemm}
\begin{proof}
There is a (elementwise least) increasing sequence $\{c_k\}_{k>0}$ such that for each $k_0>0$ and $k\geq k_0$ we have $\pPATH_k^{k-{c_{k_0}}}\neq \varnothing$, moreover $\lim_{k\to\infty} c_k = \infty$.

$\pPATH^{k-c}_k$ are disjoint for different choices of $c$, so by direct computation we have that the fraction $\abs{\G_k'}/\abs{\pPATHl_k}$ is
\begin{align}
\sum_{c=0}^{c_k} \frac{\abs{\pPATH_k^{k-c}}}{\abs{\pPATHl_k}}&=\sum_{c=0}^{c_k}\frac{1}{c!\sum_{i=0}^{k-2}\frac{1}{i!}}\\
&=\left(\frac{1}{\sum_{i=0}^{k-2}\frac{1}{i!}}\right)\sum_{c=0}^{c_k}\frac{1}{c!},
\end{align}
which tends to $1$ as $k\to\infty$.
\end{proof}

\begin{thrm}
\[\lim_{F_{nbtree}}\pPATHl_n\bbl(\exists v)(\exists u)(\forall w)((v\not=0)\land (E(v,u)) \land(E(v,w) \to u=w))\bbr=\0\]
\end{thrm}
\begin{proof}(sketch)
The probability that a tree fails on any of the wide sequences $\pPATH_k^{k-c}$ is infinitesimal due to Lemma \ref{lemmpathnnc}. This is true for any of their unions so for $\G_k'$ in particular. In Lemma \ref{lemmpathmaj} we proved that $\G_k'$ is a major portion of $\pPATHl_k$, so by Lemma \ref{lemmporup} the probability that a tree fails in $\pPATHl$ is arbitrarily close to $1$.
\end{proof}

So we have for a larger set of instances that their wide limit has no solution relative to $F_{nbtree}$. Moreover, we have the following.

\begin{crll}\label{crllpPATHl}
\[\Th\left(\lim_{F_{nbtree}}\pPATHl_n\right)=\Th\left(\lim_{F_{nbtree}}\pPATH_n\right)\]
\end{crll}
\begin{proof}
Since the wide sequences $\pPATH_k^{k-c}$, for $c\in\N$, partition a major portion of $\pPATHl_k$, and in each of them the probability of a tree in $F_{nbtree}$ inspecting a degree 0 vertex is infinitesimal, we have that the wide limit of the major portion $\G_k'$ satisfies the theory $\Th(\lim_{F_{nbtree}}\pPATH_n)$ and therefore so does $\lim_{F_{nbtree}}\pPATHl_n$ .
\end{proof}



\section{\texorpdfstring{$\G_k=\pDPATH_k$}{Gk=*DPATHk}}\label{secpDPATH}

As $\pPATH_k$ was the wide sequence consisting of the hardest instances of $\LEAF$ the complete problem for $\PPA$, we define $\pDPATH_k$ analogously but in the directed case so it consists of the hardest instances of $\SOURCEORSINK$ the complete problem for $\PPAD$.


\begin{defi}
We define $\pDPATH_k$ (the pointed directed paths on $k$ vertices) as the set of all directed graphs $G$ on the vertex set $\set{k}$, where $G$ is isomorphic to the path on $k$ vertices such that $\deg_G^{+}(0)=0$ and $\deg_G^{-}(1)=1$.
\end{defi}

But now, since we are working with directed graphs which have two types of neighbour sets $N_G^+(v)=\{w\in V_G;E_G(w,v)\}$ and $N_G^-(v)=\{w\in V_G;E_G(v,w)\}$, we would like to define a family $F_{dtree}$ of those trees which can inspect either of the neighbour sets.

\begin{defi}
After we fix $n$, we define $F_{dtree}$ as the set of all functions computed by some labeled tree with the following shape:

\begin{itemize}
\item Each non-leaf node is labeled by some $v\in\set{n}$ and a symbol $\circ\in\{+,-\}$. 
\item For each $v \in \set{n}$ and a node $w$ there is an outgoing edge from $w$ labeled $\{v\}$ and also an outgoing edge labeled $\varnothing$.
\item Each leaf is labeled by some $m\in \set{g_n}$.
\item The depth of the tree is at most $g_n^{1/t}$ for some $t>\mathbb{N}$.
\end{itemize}

Computation of such a tree on an undirected graph $G$ goes as follows. We interpret the non-leaf nodes as questions "what is $N_G^\circ(v)$?" and the edges as answers from our graph $G$, and thus we follow a path determined by $G$ until we find a leaf, in which case the computation returns the label of the leaf.

We denote the set of such trees $\T_{dtree}$.
\end{defi}

We will not cover details, but analysis of these trees in $\T_{dtree}$ finding the nonzero sink is more or less the same as the $F_{nbtree}$ case for $\pPATH_k$, so we have the following.

\begin{thrm}
\[\lim_{F_{nbtree}}\pDPATH_n\bbl(\exists v)(\forall w)((v\not=0)\land \lnot E(v,w))\bbr=\0\]
\end{thrm}

\begin{crll}\label{crllpDPATH}
$\Th(\lim_{F_{nbtree}}\pDPATH_n)$ is complete.
\end{crll}

In the type 2 complexity theory of $\TFNP^2$ we know that there is no (oracle polynomial time) reduction from $\LEAF$ to $\SOURCEORSINK$. An important question arises -- is this reflected in the second order arithmetical expansion $K(\pPATH_n,F_{nbtree},G_{nbtree})$? Where $G_{nbtree}$ is defined analogously as $G_{rud}$ but with the components in the tuple from $F_{nbtree}$. More specifically we ask the following.

\begin{ques}
Consider an instance of $\SOURCEORSINK$ defined by some $\Theta\in G_{nbtree}$. Let $\varphi_\Theta$ be the first order statement which says `$\Theta$ has a solution'. Are all $\varphi_\Theta$ valid in $K(\pPATH_n,F_{nbtree},G_{nbtree})$?
\end{ques}


\begin{defi}
We say that $\{\mathcal{G}_k\}_{k=0}^\infty$ is \textbf{isomorphism closed}, if there is $k_0$ such that for every $k>k_0$ if we have that $G_1\in\mathcal{G}_k$, $V_{G_2}=[g_k]$ and $G_1\cong G_2$ then $G_2\in\mathcal{G}_k$.

We say that $\{\mathcal{G}_k\}_{k=0}^\infty$ is \textbf{categorical} if there is $k_0$ such that for every $k>k_0$ if we have $G_1,G_2\in\mathcal{G}_k$ then $G_1\cong G_2$. For a categorical sequence $\{G_k\}_{k=0}^\infty$ we denote $G_k$ the lexicographically minimal element of $\mathcal{G}_k$.
\end{defi}

\begin{lemm}
Let $\{\mathcal{G}_k\}_{k=0}^\infty$ be categorical and isomorphism closed, then for large enough $k$
\[\abs{\mathcal{G}_k}=\frac{g_k!}{\abs{\Aut(G_k)}}.\]
\end{lemm}
\begin{proof}
Every $\rho\in S_{g_k}$ defines an isomorphism $\rho:G_k\to\rho(G_k)$, where $\rho(G_k)$ is a graph obtained from $G_k$ by renaming every vertex $v$ to $\rho(v)$.

\textbf{Claim:} For any $\rho,\pi\in S_{g_k}$:
\[\rho(G_k)=\pi(G_k)\iff\exists \tau\in\Aut(G_k):\rho\circ\tau=\pi.\]
\textit{Proof of claim.} ``$\Rightarrow$" Let $\rho(G_k)=\pi(G_k)$, therefore $\tau:=\rho^{-1}\circ\pi\in\Aut(G_k)$ and $\rho\circ\tau=\rho\circ\rho^{-1}\circ\pi=\pi.$

``$\Leftarrow$" Let $\rho\circ\tau=\pi$. Then $\pi(G_k)=\rho(\tau(G_k))=\rho(G_k). \qedhere$

Notice that the $\tau$ in the statement of the claim is uniquely determinted by $\rho^{-1}\circ\pi$. Therefore if we defined a quotient set $S_{g_k}/\sim$ with $\rho\sim\pi\iff\rho(G_k)=\pi(G_k)$ then $\abs{S_{g_k}/\sim}=\frac{g_k!}{\abs{\Aut(G_k)}}$.

The Lemma follows from noticing that if we start with $\{G_k\}$ and then we build $\mathcal{G}_k$ by finding isomorphic graphs on the vertex set $[g_k]$ we can only do so by trying different permutation from $S_{g_k}$ and these permutations find the same graph if and only if they are in the same $\sim$-class. Therefore there is a bijection between $S_{g_k}/\sim$ and $\mathcal{G}_k$. \qed
\end{proof}

\begin{lemm}[Candidate for optimal search trees]
Let $\{\mathcal{G}_k\}_{k=0}^\infty$ be categorical and isomorphism closed, let $\varphi(x_0,\dots,x_{l-1})$ be an open $\{E\}$-formula, let $\models \varphi(\overline x) \to \Land_{i,j=0}^{l-1,l-1}x_i=^{b_{ij}}x_j$ for some $b_{ij}\in\{0,1\}$, let $k_0\geq0$ and define $\{q_k\}_{k=k_0}^\infty$ as follows

\[q_k:=\frac{g_k!}{\abs{\Aut(G_k)}}\cdot\frac{\abs{\varphi(G_k)}}{\abs{\bigcup_{G\in\mathcal{G}}\varphi(G_k)}}.\]

Then there is $c\in\mathbb{N}$ and trees $T_0,\dots,T_{l-1}$ of depth $n^{(r)} \cdot c$, (with $n^{(r)}$ being defined in the proof) such that for the $\overline \alpha$ computed by $\overline T$ we have $\bbl\varphi(\overline \alpha)\bbr=\1$.
\end{lemm}

\begin{proof}
We will use the identity from the statement to construct a search tree (iterated $T_{\overline a}$) which almost always finds a witness to $\varphi$.

We will analyze the problem in the finite case for big enough $k>0$. We should only check those tuples included in $\bigcup_{G\in\mathcal{G}_k}\varphi(G)$. For example, if we are trying to find an edge then we need not check the constant tuples $(a,a)$. Moreover, to succeed we only need to check one specific tuple in each $\varphi(G),G\in\mathcal{G}_k$.

Consider the set $S=\{(G,\overline a); G\in \mathcal{G}_k, G\models\varphi(\overline a)\}$ and a projection to the second coordinate $p_2: S\to\bigcup_{G\in\mathcal{G}_k}\varphi(G)$. Since $\abs{S}=\frac{g_k!}{\abs{\Aut(G)}}\cdot\abs{\varphi(G_k)}$ we have that $q_k$ is the average size of a $p_2$ preimage of any $\overline a\in \bigcup_{G\in\mathcal{G}_k}\varphi(G)$. 

\textbf{Claim:} For all $\overline a,\overline b\in\bigcup_{G\in\mathcal{G}_k}\varphi(G)$ we have $\abs{p_2^{-1}[\overline a]}=\abs{p_2^{-1}[\overline b]}=q_k$.

\textit{Proof of claim.} We will prove that for any $\overline a,\overline b\in\bigcup_{G\in\mathcal{G}_k}\varphi(G)$ we have $\abs{p_2^{-1}[\overline a]}\leq\abs{p_2^{-1}[\overline b]}$, by symmetry, they must be equal and also equal to $q_k$ which is the average size of any singleton preimage.

Let $p_2^{-1}[\overline a]=\{G_0,\dots,G_{s-1}\}\times\{\overline a\}$ and let $\rho=(b_0\: a_0)\dots(b_{l-1}\:a_{l-1})$, this is a permutation from the condition on $\varphi$. Then \[p_2^{-1}[\overline b]\supseteq\{\rho(G_0),\dots,\rho(G_{s-1})\}\times\{\rho(\overline a)=\overline b\}.\qed\] 

Now consider the multiset $M=(\bigcup_{G\in\mathcal{G}_k}\varphi(G),\text{count}:\overline a \mapsto \abs{p_2^{-1}[\overline a]})$, we will construct the searching tree by plucking elements from this multiset in the following way.

Let $M^{(0)}:=M,\mathcal{G}_k^{(0)}=\mathcal{G}_k$. For $i\geq0$ and $M^{(i)},\mathcal{G}_k^{(i)}$ built, take some $\overline a\in M^{(i)}$ with maximal $\text{count}(\overline a)$, put $\mathcal{G}_k^{(i+1)}=\mathcal{G}_k^{(i)}\setminus p_2^{-1}[\overline a]$ and form $M^{(i+1)}$ by removing $\overline a$, and for every $\overline b \in p_1[p_2^{-1}[\overline a]]\setminus\{\overline a\}$ setting $\text{count}_{M^{(i+1)}}(\overline b)=\max\{0,\text{count}_{M^{(i)}}(\overline b)-(\varphi(G_k))\}$. We also add $T_{\overline a}$ to the leaves of the tree we are constructing $T_i$ and call it $T_{i+1}$.

For each $i\geq 0$ we have that $T_i$ finds a witness in $G\in\mathcal{G}_k$ iff $G\not\in\mathcal{G}_k^{(i)}$. So to calculate the probability of success of $T_i$ we just need to find upper bounds on the cardinality of $\mathcal{G}_k^{(i)}$. 

Define $m_i:=\max\{\text{count}(\overline a);\overline a \in M^{(i)}\}$. Let $k^{(0)}\geq 0$ be the greatest number such that for all $i<k^{(0)}$: 
\(m_i=q_k.\) 

Define a set $M_m^{(i)}=\{\overline a;\text{count}_{M_i}(\overline a)=m_i\}$. We can see, that $k^{(0)}\geq1$ and $M^{(0)}=\bigcup_{G\in\mathcal{G}_k}\varphi(G)$. At each step $i<k^{(0)}$ we construct $T_{i+1}$ by searching for some $\overline a \in M_{m}^{(i)}$, this results in $\abs{\mathcal{G}_k^{(i+1)}}=\abs{\mathcal{G}_k^{i}}-q_k$. We also remove one instance of every $\overline b \in p_1[p_2^{-1}(\overline a)]\setminus\{\overline a\}$ from $M^{(i)}$ to form $M^{(i+1)}$, this results in $\abs{M^{(i+1)}_m}\geq \abs{M^{(i)}_m}-1-q_k\cdot(\abs{\varphi(G_k)-1}).$

Therefore 

\begin{align}
k^{(0)}&\geq \left \lfloor \frac{\abs{M^{(0)}_m}}{q_k\cdot\abs{\varphi(G_k)}} \right \rfloor\\
&=
\left\lfloor\frac{\abs{\bigcup_{G\in\mathcal{G}_k}\varphi(G)}}{q_k\cdot\abs{\varphi(G_k)}}\right\rfloor
,
\end{align}

and $\abs{\mathcal{G}_k^{(k^{(0)})}}=\abs{\mathcal{G}_k}-k^{(0)}\cdot q_k=\frac{\abs{\Aut(G_k)}}{g_k!}-\left\lfloor\frac{\abs{\cup_{G\in\mathcal{G}_k}\varphi(G)}}{q_k\cdot\abs{\varphi(G_k)}}\right\rfloor\cdot q_k\leq \frac{\abs{\Aut(G_k)}}{g_k!}-\left\lfloor\frac{\abs{\cup_{G\in\mathcal{G}_k}\varphi(G)}}{\abs{\varphi(G_k)}} \right \rfloor.$

However the right hand side of the last inequality is rarely $\leq 0$, so generally one has to continue with plucking even after $k^{(0)}$-many steps. We define $k^{(j)}$ as the greatest number such that for all $i<k^{(j)}:m_i\geq q_k-j$ and continue for $k^{(r)}$ steps, where $r$ is the smallest number such that


\begin{align}
\abs{\mathcal{G}_{k}^{k^{(r)}}}&=\abs{\mathcal{G}_k}-k^{(0)}\cdot q_k- \sum_{j=1}^{r}(k^{(j)}-k^{(j-1)})\cdot (q_k - j)\\
&=0.
\end{align}

However, this requires a general analysis of $k^{(j)}$ and I haven't manage to compute that.

For $k=n$ in $\mathcal{M}$ we put $n^{(r)}:=k^{(r)}.$
\end{proof}


\chapter*{Concluding remarks}
\addcontentsline{toc}{chapter}{Concluding remarks}

In this thesis we built a basic theory around wide limits of graphs, proved several general theorems and described the theories of the wide limits $\lim_F\EDGE_n$, $\lim_{F_{rud}}\ALL_n$, $\lim_{F_{nbtree}}\pPATH_n$, $\lim_{F_{nbtree}}\pPATHl_n$ and $\lim_{F_{dtree}}\pDPATH_n$ and proved that they are complete (Corollary \ref{crlltoosparse}, Theorem \ref{thrmall} and corollaries \ref{crllpPATH}, \ref{crllpPATHl} and \ref{crllpDPATH}).  We also proved that $\lim_{F_{rud}}^{G_{rud}}\SK_n^{1/2}$ contains a clique of size $\floor{n/(2\ln n)}$ (Theorem \ref{thrmsk}) and also that nonexistence of finite cliques in not valid in the more complex wide limit $\lim_{F_{rud}}^{G_{rud}}\CK_n^{1/2}$ (Theorem \ref{thrmck}).

During the development we planned to analyze wide limits the family $F_{poly}$ of polynomial functions. In \cite{krajicek2010forcing} it was proven that forcing with $F_{poly}$ results in quantifier elimination which implies that if an $\{E\}$-sentence holds in large enough $\G_k$, it has to hold in the limit. However the second order limit can still provide some information about the ability of polynomial time functions to search interesting subsets of $G\in\G_k$. In the end we did not get any new results about it. We want to mention that even though it would seem that the properties of $F_{poly}$ limits directly correspond to the $\P$ vs. $\NP$ question it seems that something a bit different happens. The way the limit objects are defined, it is not enough that some polynomial time algorithm does not exist to see that we cannot witness some property in the limit, but it is important that no polynomial algorithm succeeds on any standard fraction of inputs. This more closely resembles the generic case polynomial time (see \cite{gilman2007report}).

Another natural question would be to consider structures over general languages than just the language of graphs, and allow trees to query atomic sentences. Other combinatorial structures like hypergraphs and tournaments could be considered. Furthermore, wide limits of finite universal algebras could be considered which could require a whole new theory. This all leads to the fact that generalized spectra, elementary classes of $\Sigma_1^1$-logic (see \cite{Fagin74}), with restricted vertex sets make up a wide sequence, there could be a connection to the theory of spectra of sentences.


Considering structures over arbitrary vocabularies opens the door to connecting the concept of a wide limit with the complexity of the Constraint Satisfaction Problem. Solutions to $\text{CSP}(\A)$ form a wide sequence and depending on the tractability of the class we could expect different behaviour from $\lim_{F}\text{CSP}(\A)_n$.

Finally, important direction in studying wide sequences would be to characterize some $\lim_F\G_n$ without the direct construction and therefore to prove upper and lower bounds for search problems in $F$.


%%% Bibliography
%%% Bibliography (literature used as a source)
%%%
%%% We employ bibTeX to construct the bibliography. It processes
%%% citations in the text (e.g., the \cite{...} macro) and looks up
%%% relevant entries in the bibliography.bib file.
%%%
%%% The \bibliographystyle command selects, which style will be used
%%% for references from the text. The argument in curly brackets is
%%% the name of the corresponding style file (*.bst). Both styles
%%% mentioned in this template are included in LaTeX distributions.

\bibliographystyle{plainnat}    %% Author (year)
% \bibliographystyle{unsrt}     %% [number]

\renewcommand{\bibname}{Bibliography}

%%% Generate the bibliography. Beware that if you cited no works,
%%% the empty list will be omitted completely.

\bibliography{bibliography}

%%% If case you prefer to write the bibliography manually (without bibTeX),
%%% you can use the following. Please follow the ISO 690 standard and
%%% citation conventions of your field of research.

% \begin{thebibliography}{99}
%
% \bibitem{lamport94}
%   {\sc Lamport,} Leslie.
%   \emph{\LaTeX: A Document Preparation System}.
%   2nd edition.
%   Massachusetts: Addison Wesley, 1994.
%   ISBN 0-201-52983-1.
%
% \end{thebibliography}


%%% Figures used in the thesis (consider if this is needed)
%%\listoffigures

%%% Tables used in the thesis (consider if this is needed)
%%% In mathematical theses, it could be better to move the list of tables to the beginning of the thesis.
%%\listoftables

%%% Abbreviations used in the thesis, if any, including their explanation
%%% In mathematical theses, it could be better to move the list of abbreviations to the beginning of the thesis.
%%\chapwithtoc{List of Abbreviations}

%%% Attachments to the master thesis, if any. Each attachment must be
%%% referred to at least once from the text of the thesis. Attachments
%%% are numbered.
%%%
%%% The printed version should preferably contain attachments, which can be
%%% read (additional tables and charts, supplementary text, examples of
%%% program output, etc.). The electronic version is more suited for attachments
%%% which will likely be used in an electronic form rather than read (program
%%% source code, data files, interactive charts, etc.). Electronic attachments
%%% should be uploaded to SIS and optionally also included in the thesis on a~CD/DVD.
%%% Allowed file formats are specified in provision of the rector no. 72/2017.
%% \appendix
%% \chapter{Attachments}

%% \section{First Attachment}

%% \openright
\end{document}
